\documentclass[11pt,a4paper]{article}
\usepackage[T1]{fontenc}
\usepackage[utf8]{inputenc}
\usepackage{lmodern}
\usepackage{amsmath,amssymb,amsthm,mathtools}
\usepackage[margin=24mm,headheight=14pt]{geometry}
\usepackage{microtype}
\usepackage{xcolor}
\usepackage{booktabs}
\usepackage{enumitem}
\usepackage{fancyhdr}
\usepackage[colorlinks=true,linkcolor=blue!45!black,urlcolor=blue!45!black]{hyperref}
\hypersetup{pdftitle={Polynomial Rectification of the Shastri Trefoil and a Harmonic Figure-Eight},pdfauthor={Ivan Dorofeev and Alexander Frolov},pdfsubject={Exact polynomial constructions and a computational certificate}}
\pagestyle{fancy}
\fancyhf{}
\fancyhead[L]{\small Polynomial rectification}
\fancyhead[R]{\small I. Dorofeev and A. Frolov}
\fancyfoot[C]{\thepage}
\renewcommand{\headrulewidth}{0.3pt}
\setlength{\parindent}{0pt}
\setlength{\parskip}{5pt}
\setlength{\emergencystretch}{2em}
\setlist{nosep,leftmargin=*}
\newtheorem{theorem}{Theorem}[section]
\newtheorem{proposition}[theorem]{Proposition}
\newcommand{\A}{\mathbb A}
\newcommand{\Q}{\mathbb Q}
\newcommand{\C}{\mathbb C}
\newcommand{\id}{\mathrm{id}}
\newcommand{\comp}{\circ}
\title{\vspace{-15mm}\LARGE\bfseries Polynomial Rectification of the Shastri Trefoil\newline and a Harmonic Figure-Eight}
\author{Ivan Dorofeev\\[3pt]\normalsize Independent Researcher\\\normalsize\href{mailto:qft98romanov@gmail.com}{\texttt{qft98romanov@gmail.com}}\and Alexander Frolov\\[3pt]\normalsize Researcher\\\small Higher School of Modern Mathematics,\\\small MIPT}
\date{24 October 2026}

\begin{document}
\maketitle
\thispagestyle{empty}
\vspace{-4mm}

\begin{abstract}
We give an explicit implementation of coordinate transfer through affine
modifications, including repeated exceptional factors. We then
construct polynomial automorphisms rectifying the Shastri trefoil and a
harmonic figure-eight over $\Q(\sqrt{-3})$, with explicit polynomial inverses.
\end{abstract}

\section*{Acknowledgements}
The authors express their sincere gratitude to Professor Alexei Kanel-Belov
and Professor Alexander Gasnikov.


\smallskip
\begin{quote}\small
The results were obtained on 24 October 2026 through interactive work with
GPT Astra. The main polynomial identities and both explicit rectifications
have been checked by exact symbolic computation. A reproducible
computational certificate is provided in \texttt{verify.py}, with its run
log in \texttt{verification.log}.
The coordinate-only formulas were corrected on 25 October 2026. All twelve displayed components are additionally checked against the compact construction by \texttt{verify\_printed\_formulas.py}; see the accompanying errata.
\end{quote}

\section{Methods}

\subsection{The embedding problem and its sources}

The embedding problem asks whether every closed embedding
$\gamma:\A^1_{\C}\hookrightarrow\A^3_{\C}$ admits
$\Phi\in\operatorname{Aut}_{\C}(\A^3)$ with
$\Phi(\gamma(t))=(t,0,0)$. The plane case is the theorem of
Abhyankar--Moh and Suzuki \cite{AM,Suzuki}. The spatial question has a
different history: Abhyankar's Kyoto 1977 conjectures are recorded in
\cite[Chapter 11, pp.~413--414]{Abhyankar}, as documented in
\cite[Introduction, pp.~115--116]{Craighero} and
\cite[\S1, p.~11]{Shastri}. In particular, the historical conjecture
asserted the existence of nonrectifiable spatial lines.

Shastri's trefoil is taken from \cite[\S3, Example~1]{Shastri}.
The harmonic figure-eight comes from the Chebyshev-knot literature
\cite[Introduction, Definition~1.2 and the following paragraph]{KP}.
The coordinate-lift strategy has a precedent in
\cite[Proposition~3, pp.~117--118]{Craighero}: starting with a polynomial
restricting to $t$, one changes it by an element of the curve ideal.
Our output consists of explicit automorphisms, inverses and certificates
for the two specified parametrizations.

\subsection{Exact division and global polynomial identities}

Work over a field $k$ of characteristic zero. A polynomial map is called an
automorphism when a polynomial inverse is supplied. The basic transition is
\[
\boxed{\text{invertible after localization}
\; +\; \text{polynomial in both directions}
\;\Longrightarrow\; \text{polynomial automorphism}.}
\]
Indeed, for $0\ne m\in k[x]$, the inclusion
\[
k[x,y,z]\hookrightarrow k[x,1/m,y,z]
\]
is injective. Once both maps have polynomial components, composition
identities proved after adjoining $1/m$ hold in the original polynomial
ring. Thus the fibres $m(x)=0$ are included.

\subsection{Transfer of coordinates}

Suppose polynomial inverses are given:
\begin{equation}\label{eq:given}
\begin{aligned}
A(x,y,z)&=(x,H(x,y,z),W(x,y,z)),\\
A^{-1}(x,T,V)&=(x,R(x,T,V),S(x,T,V)).
\end{aligned}
\end{equation}
Choose $0\ne m(x)\in k[x]$ and $p(x,y)\in k[x,y]$. The modification
\[
\beta_m(x,y,z)=(x,p(x,y)+m(x)z,y)
\]
is invertible after adjoining $1/m$. We seek a correction that makes the
complete transfer and its inverse polynomial.

\begin{theorem}[Sufficient transfer criterion]\label{thm:transfer}
Assume $\deg m>0$ and put $B=k[x]/(m)$. Write
\[
c(y)=H(x,p(x,y),y)=c_0+c_1y+\sum_{j=2}^{d}c_jy^j
\quad\text{in }B[y].
\]
If $c_1$ is a unit in $B$ and every $c_j$, $j\ge2$, is nilpotent, then
$H(x,p(x,y)+m(x)z,y)$ extends to a polynomial automorphism preserving $x$.
Both the extension and its inverse are obtained by finite polynomial
operations.
\end{theorem}

The existence conclusion is a special case of the residual-variable
theorem of Bhatwadekar--Dutta \cite[Theorem~3.2]{BD}; the constructive relation to
affine modifications is explained below. We give the formulas and an
elementary proof in order to make the companions and inverses explicit.

Over an algebraic closure, the hypothesis is equivalent to
\[
\deg_y H(r,p(r,y),y)=1\qquad\text{for every root }r\text{ of }m.
\]
The construction below retains the full multiplicities in $m$.

\paragraph{Step 1: invert the centre modulo $m$.}
Let $J=(c_2,\ldots,c_d)\subset B$ and choose $e\ge1$ with $J^e=0$.
One may take the largest exponent in the irreducible factorization of $m$.
In $B[T]$, set
\begin{equation}\label{eq:iteration}
\begin{aligned}
\psi_0(T)&=c_1^{-1}(T-c_0),\\
\psi_{j+1}(T)&=c_1^{-1}
\left(T-c_0-\sum_{\ell=2}^{d}c_\ell\psi_j(T)^\ell\right),
\qquad 0\le j<e.
\end{aligned}
\end{equation}
The differences satisfy $\psi_{j+1}-\psi_j\in J^{j+1}[T]$: subtract
successive iterations and factor every difference of powers. Consequently
the iteration stabilizes, and
\begin{equation}\label{eq:center-inverse}
c(\psi_e(T))=T,\qquad \psi_e(c(y))=y.
\end{equation}
For the second identity, $c(v)-c(y)$ is $v-y$ times $c_1+n$, with
$n\in J B[y,v]$. After substituting $v=\psi_e(c(y))$, this factor is a unit;
its inverse is a finite geometric sum. Cancel it to obtain the identity.

\paragraph{Step 2: choose the correction.}
Take any representative $\varphi(x,T)\in k[x,T]$ of
\begin{equation}\label{eq:correction}
\varphi(x,T)\equiv
W\bigl(x,p(x,\psi_e(T)),\psi_e(T)\bigr)\pmod{m(x)}.
\end{equation}
For example, choose the remainder of degree less than $\deg m$ in $x$.

\paragraph{Step 3: form the maps.}
Define $\mathcal A$ componentwise by
\begin{equation}\label{eq:transfer}
\begin{aligned}
\mathcal A_1(x,y,z)&=x,\\
\mathcal A_2(x,y,z)&=H(x,p(x,y)+m(x)z,y),\\
\mathcal A_3(x,y,z)&=
\frac{\begin{gathered}
W(x,p(x,y)+m(x)z,y)\\[-2pt]
-\varphi\bigl(x,H(x,p(x,y)+m(x)z,y)\bigr)
\end{gathered}}{m(x)}.
\end{aligned}
\end{equation}
Its inverse is
\begin{equation}\label{eq:transfer-inverse}
\begin{aligned}
(\mathcal A^{-1})_1(x,T,Q)&=x,\\
(\mathcal A^{-1})_2(x,T,Q)&=S(x,T,m(x)Q+\varphi(x,T)),\\
(\mathcal A^{-1})_3(x,T,Q)&=
\frac{\begin{gathered}
R(x,T,m(x)Q+\varphi(x,T))\\[-2pt]
-p\bigl(x,S(x,T,m(x)Q+\varphi(x,T))\bigr)
\end{gathered}}{m(x)}.
\end{aligned}
\end{equation}
Every displayed quotient is exact polynomial division of the entire numerator.

\paragraph{Step 4: prove both divisions exact.}
Equations \eqref{eq:center-inverse}--\eqref{eq:correction} give
\[
\varphi(x,H(x,p(x,y),y))\equiv W(x,p(x,y),y)\pmod m.
\]
Replacing $p$ by $p+mz$ preserves this congruence, proving the divisibility
in \eqref{eq:transfer}.

Apply $A^{-1}$ to the two centre identities. In $B[T]$ this gives
\[
\begin{aligned}
R(x,T,\varphi(x,T))&=p(x,\psi_e(T)),\\
S(x,T,\varphi(x,T))&=\psi_e(T).
\end{aligned}
\]
Replacing $\varphi$ by $mQ+\varphi$ preserves the congruences. Their
difference proves divisibility of the numerator in
\eqref{eq:transfer-inverse} by the full $m$.

\paragraph{Step 5: conclude global invertibility.}
Over $k[x,1/m]$, let
\[
C_m(x,T,V)=\left(x,T,\frac{V-\varphi(x,T)}{m(x)}\right),\qquad
C_m^{-1}(x,T,Q)=(x,T,m(x)Q+\varphi(x,T)).
\]
Then
\[
\mathcal A=C_m\comp A\comp\beta_m,\qquad
\mathcal A^{-1}=\beta_m^{-1}\comp A^{-1}\comp C_m^{-1}.
\]
Both compositions are identities after localization. Polynomiality and
injectivity now give
\[
\mathcal A^{-1}\comp\mathcal A=\id,
\qquad \mathcal A\comp\mathcal A^{-1}=\id
\]
globally. This proves Theorem~\ref{thm:transfer}.
If $m$ is a nonzero constant, take $\varphi=0$; both divisions are by a
scalar and the same formulas apply. Changing $\varphi$ to $\varphi+mD(x,T)$
only changes the third coordinate by the polynomial shear
$-D(x,\mathcal A_2)$.

\subsection{Relation to established constructions}

The map $\beta_m$ is the affine modification with divisor $m=0$ and
centre ideal $(m,Y-p(x,Z))$. Under $A=(x,H,W)$, inversion of
$H(x,p(x,Z),Z)$ modulo the full $m$ identifies this centre with
$(m,V-\varphi(x,T))$. Thus the transferred map is the lift between these
two modifications furnished by \cite[Proposition~2.1]{KZ};
\cite[Corollary~2.2]{KZ} gives the centre-preserving case.

The quotient-ring step uses the classical description of one-variable
automorphisms over a ring: a unit linear coefficient and nilpotent higher
coefficients \cite[Proposition~1]{EV}. Related finite constructions and
transfer arguments occur in \cite[Theorem~4 and \S5]{EV}.
A particularly close formula is recalled in \cite[\S1, Lemma~1.1]{EVM}:
if $b\circ c=c\circ b=\id\pmod m$, then
\[
F(u,v)=\left(mv+c(u),\frac{b(mv+c(u))-u}{m}\right)
\]
is an automorphism, with inverse
\[
F^{-1}(U,V)=\left(b(U)-mV,\frac{U-c(b(U)-mV)}m\right).
\]
Both divisions are exact, and cancellation proves the inverse identities.

There is also an existence proof using \cite[Theorem~3.2]{BD}. Over the base $k[x]$,
$H\circ\beta_m$ is a coordinate on the generic fibre and on every fibre
where $m$ is nonzero. On a fibre dividing $m$, the hypothesis makes it
linear with nonzero slope in $y$. It is therefore a residual coordinate
in two variables, so the theorem of Bhatwadekar--Dutta applies.
The proof above supplies explicit formulas without invoking that theorem.
No novelty claim is made for the general lifting principle; the concrete
rectifiers and their certificates are the focus here.

\subsection{From a relative coordinate to a straight line}

Suppose \eqref{eq:given} holds and
$\gamma(t)=(b(t),c(t),d(t))$ satisfies $H(\gamma(t))=t$.
Put $w(t)=W(\gamma(t))$. The rectifier is
\begin{equation}\label{eq:straighten}
\Phi(x,y,z)=\bigl(H(x,y,z),\ x-b(H(x,y,z)),\
W(x,y,z)-w(H(x,y,z))\bigr).
\end{equation}
Its inverse is explicitly
\begin{equation}\label{eq:straighten-inverse}
\begin{aligned}
(\Phi^{-1})_1(X,Y,Z)&=Y+b(X),\\
(\Phi^{-1})_2(X,Y,Z)&=R(Y+b(X),X,Z+w(X)),\\
(\Phi^{-1})_3(X,Y,Z)&=S(Y+b(X),X,Z+w(X)).
\end{aligned}
\end{equation}
These maps compose $A$ with a permutation and two polynomial shears.
Thus $\Phi(\gamma(t))=(t,0,0)$. The input is an already verified
coordinate pair $(H,W)$, not merely a polynomial lifting the parameter.

\clearpage
\section{The Shastri trefoil}

We work over
\[
K=\Q[a]/(a^2-a+1)\simeq\Q(\sqrt{-3}),\qquad
a=\frac{1+i\sqrt3}{2},\qquad a(1-a)=1.
\]
The input curve, given in \cite[\S3, Example~1, pp.~14--15]{Shastri}, is
\begin{equation}\label{eq:trefoil}
\gamma_T(t)=(t^3-3t,\ t^4-4t^2,\ t^5-10t).
\end{equation}
Shastri already gives the parameter lift
$yz-x^3-5xy+2z-7x$; the construction below supplies a coordinate lift
together with a complete automorphism and its inverse.

\subsection{The base coordinate pair and its inverse}

Define
\begin{equation}\label{eq:H}
\begin{aligned}
H(x,y,z)={}&ax^3y^2+2x^2yz+(1-a)xz^2\\
&+\bigl((a-1)x^3-ax\bigr)y+(-1+ax^2-ax^4)z\\
&-2x^5+(6-2a)x^3-2(1-a)x,
\end{aligned}
\end{equation}
and
\begin{equation}\label{eq:f}
\begin{aligned}
f(x,T)={}&-T+x\bigl((1-a)T^2+2a-2\bigr)\\
&+x^2(2aT^3-5aT)+x^3(-5T^4+15T^2-2a)\\
&+x^4\bigl(14(1-a)T^5+50(a-1)T^3+(29-24a)T\bigr).
\end{aligned}
\end{equation}
The two exact quotients
\begin{equation}\label{eq:base-divisions}
\begin{aligned}
W(x,y,z)&=\frac{a\bigl(axy+z-f(x,H(x,y,z))\bigr)}{x^5},\\
\chi(x,T)&=\frac{a\bigl(T-H(x,0,f(x,T))\bigr)}{x^5}
\end{aligned}
\end{equation}
are polynomials. Reduction by $a^2-a+1$ and coefficientwise division
by $x^5$ give their full coefficients in the certificate.
Define
\begin{equation}\label{eq:RS}
\begin{aligned}
R(x,T,V)={}&-\chi(x,T)\\
&+\bigl(-1+ax^2-ax^4+2(1-a)xf(x,T)\bigr)V\\
&+(1-a)^2x^6V^2,\\[2pt]
S(x,T,V)={}&f(x,T)+(1-a)x^5V-axR(x,T,V).
\end{aligned}
\end{equation}
Then
\[
A_0=(x,H,W),\qquad A_0^{-1}=(x,R,S).
\]

\subsection{Certificate of both inverse identities}

The exact polynomial checks are
\begin{align}
x^5W(x,y,z)&=a\bigl(axy+z-f(x,H(x,y,z))\bigr),\label{eq:cert1}\\
H(x,R(x,T,V),S(x,T,V))&=T,\label{eq:cert2}\\
axR(x,T,V)+S(x,T,V)-f(x,T)&=(1-a)x^5V,\label{eq:cert3}\\
R(x,H(x,y,z),W(x,y,z))&=y,\label{eq:cert4}\\
S(x,H(x,y,z),W(x,y,z))&=z.\label{eq:cert5}
\end{align}
Substitute $R,S$ into \eqref{eq:cert1} and use
\eqref{eq:cert2}--\eqref{eq:cert3}:
\[
x^5W(x,R,S)=a(1-a)x^5V=x^5V.
\]
Cancellation in $K[x,T,V]$ gives $W(x,R,S)=V$.
Together with \eqref{eq:cert2}, this proves $A_0\comp A_0^{-1}=\id$;
\eqref{eq:cert4}--\eqref{eq:cert5} prove the opposite composition.
The divisions in \eqref{eq:base-divisions} and this cancellation are
polynomial operations, valid also on the fibre $x=0$.

\subsection{Normalization and final rectifier}

The affine maps
\[
M(x,y,z)=(x,y+2,z-5x),\qquad M^{-1}(x,y,z)=(x,y-2,z+5x)
\]
give
\[
M(\gamma_T(t))=(t^3-3t,\ t^4-4t^2+2,\ t^5-5t^3+5t).
\]
Put
\[
h_T(x,y,z)=H(x,y+2,z-5x),\qquad
w_T(x,y,z)=W(x,y+2,z-5x).
\]
Exact substitution yields
\[
h_T(\gamma_T(t))=t,\qquad w_T(\gamma_T(t))=8-14t^2.
\]
The final components are therefore
\begin{equation}\label{eq:PhiT}
\begin{aligned}
(\Phi_T)_1(x,y,z)&=h_T(x,y,z),\\
(\Phi_T)_2(x,y,z)&=x-h_T(x,y,z)^3+3h_T(x,y,z),\\
(\Phi_T)_3(x,y,z)&=w_T(x,y,z)-8+14h_T(x,y,z)^2.
\end{aligned}
\end{equation}
The inverse, on input $(X,Y,Z)$, is
\begin{equation}\label{eq:PhiTinv}
\begin{aligned}
(\Phi_T^{-1})_1(X,Y,Z)&=X^3-3X+Y,\\
(\Phi_T^{-1})_2(X,Y,Z)&=R(X^3-3X+Y,X,Z+8-14X^2)-2,\\
(\Phi_T^{-1})_3(X,Y,Z)&=S(X^3-3X+Y,X,Z+8-14X^2)\\
&\quad+5(X^3-3X+Y).
\end{aligned}
\end{equation}
Explicitly, with
\[
L_T(x,T,V)=(T,x-T^3+3T,V-8+14T^2),
\]
the transition is
\[
\boxed{\Phi_T=L_T\comp A_0\comp M,\qquad
\Phi_T^{-1}=M^{-1}\comp A_0^{-1}\comp L_T^{-1}.}
\]
Each factor is a global polynomial automorphism, and
\begin{equation}\label{eq:trefoil-result}
\boxed{\Phi_T(\gamma_T(t))=(t,0,0),\qquad
\deg\Phi_T=(5,15,24),\qquad \det D\Phi_T=1.}
\end{equation}
The degree triple records the displayed construction.

\clearpage
% Exact coordinate formulas copied from trefoil_explicit/overleaf/main.tex.
% Only left-hand-side names and surrounding presentation have changed.
\subsection{Coordinate-only formulas}
The formulas below use only the source coordinates $(x,y,z)$, the target
coordinates $(X,Y,Z)$, and the fixed scalar $a$ from the preceding construction.

\subsubsection*{Forward map}
\paragraph{Component 1.}
\begingroup
\allowdisplaybreaks[4]
\begin{align*}
(\Phi_T)_1(x,y,z)={}&-z+x\bigl((1-a)z^2-ay+3\bigr)\\
&{}+x^2z(2y+11a-6)\\
&{}+x^3\bigl(ay^2+(5a-11)y+9-26a\bigr)\\
&{}-ax^4z+(5a-2)x^5
\end{align*}
\endgroup

\paragraph{Component 2.}
\begingroup
\allowdisplaybreaks[4]
\begin{align*}
(\Phi_T)_2(x,y,z)={}&x-\Bigl[-z+x\bigl((1-a)z^2-ay+3\bigr)\\
&\qquad +x^2z(2y+11a-6)\\
&\qquad +x^3\bigl(ay^2+(5a-11)y+9-26a\bigr)\\
&\qquad -ax^4z+(5a-2)x^5\Bigr]^3\\
&{}+3\Bigl[-z+x\bigl((1-a)z^2-ay+3\bigr)\\
&\qquad +x^2z(2y+11a-6)\\
&\qquad +x^3\bigl(ay^2+(5a-11)y+9-26a\bigr)\\
&\qquad -ax^4z+(5a-2)x^5\Bigr]
\end{align*}
\endgroup

\paragraph{Component 3.}
\begingroup
\footnotesize
\allowdisplaybreaks[4]
\begin{align*}
(\Phi_T)_3(x,y,z) ={}&- y + z^{6} \left(42 a - 42\right) + z^{4} \left(175 - 175 a\right) + z^{2} \left(147 a - 148\right)\\
&{}+x\Bigl[- 252 y z^{5} + 700 y z^{3} - 296 y z - 120 a z^{7} + z^{5} \left(756 - 686 a\right) + z^{3} \left(2594 a - 2148\right)\\
&\qquad + z \left(954 - 1063 a\right)\Bigr]\\
&{}+x^{2}\Bigl[- 630 a y^{2} z^{4} + 1050 a y^{2} z^{2} + y^{2} \left(- 148 a - 1\right) + y z^{6} \left(840 - 840 a\right)\\
&\qquad + y z^{4} \left(350 a + 3430\right) + y z^{2} \left(1342 a - 7782\right) + y \left(1069 - 120 a\right) + 135 z^{8}\\
&\qquad + z^{6} \left(2520 a + 1050\right) + z^{4} \left(4646 a - 7778\right) + z^{2} \left(10277 - 14168 a\right) + 1910 a - 1536\Bigr]\\
&{}+x^{3}\Bigl[y^{3} z^{3} \left(840 - 840 a\right) + y^{3} z \left(700 a - 700\right) + 2520 y^{2} z^{5} + y^{2} z^{3} \left(7560 a - 700\right)\\
&\qquad + y^{2} z \left(- 6436 a - 1346\right) + 1080 a y z^{7} + y z^{5} \left(21420 a - 15120\right) + y z^{3} \left(- 12528 a - 18576\right)\\
&\qquad + y z \left(28316 - 7806 a\right) + z^{9} \left(70 a - 70\right) + z^{7} \left(2000 a - 5240\right)\\
&\qquad + z^{5} \left(- 21414 a - 20590\right) + z^{3} \left(52712 - 20252 a\right) + z \left(33927 a - 22907\right)\Bigr]\\
&{}+x^{4}\Bigl[630 y^{4} z^{2} - 175 y^{4} + 4200 a y^{3} z^{4} + y^{3} z^{2} \left(6860 a - 7560\right) + y^{3} \left(2144 - 2594 a\right)\\
&\qquad + y^{2} z^{6} \left(3780 a - 3780\right) + y^{2} z^{4} \left(15750 a - 53550\right) + y^{2} z^{2} \left(18792 - 46644 a\right)\\
&\qquad + y^{2} \left(10233 a + 3915\right) - 630 y z^{8} + y z^{6} \left(- 22680 a - 14000\right) + y z^{4} \left(107086 - 210036 a\right)\\
&\qquad + y z^{2} \left(97440 a + 60624\right) + y \left(11032 a - 33842\right) - 14 a z^{10} + z^{8} \left(1890 - 3290 a\right)\\
&\qquad + z^{6} \left(76600 - 52550 a\right) + z^{4} \left(99328 a + 168972\right) + z^{2} \left(68747 a - 210957\right) - 30180 a - 20159\Bigr]\\
&{}+x^{5}\Bigl[252 a y^{5} z + y^{4} z^{3} \left(4200 a - 4200\right) + y^{4} z \left(- 350 a - 3430\right) - 7560 y^{3} z^{5}\\
&\qquad + y^{3} z^{3} \left(- 50400 a - 21000\right) + y^{3} z \left(31088 - 18560 a\right) - 2520 a y^{2} z^{7}\\
&\qquad + y^{2} z^{5} \left(68040 - 110040 a\right) + y^{2} z^{3} \left(420104 - 205900 a\right) + y^{2} z \left(157992 a - 97500\right)\\
&\qquad + y z^{9} \left(140 - 140 a\right) + y z^{7} \left(26320 - 11200 a\right) + y z^{5} \left(144520 a + 315300\right)\\
&\qquad + y z^{3} \left(1072852 a - 397504\right) + y z \left(- 284634 a - 136916\right) + z^{9} \left(420 a + 350\right)\\
&\qquad + z^{7} \left(56440 a - 11780\right) + z^{5} \left(575490 a - 576910\right) + z^{3} \left(- 290692 a - 749608\right)\\
&\qquad + z \left(438422 - 151320 a\right)\Bigr]\\
&{}+x^{6}\Bigl[y^{6} \left(42 a - 42\right) - 2520 y^{5} z^{2} + y^{5} \left(70 - 756 a\right) - 9450 a y^{4} z^{4}\\
&\qquad + y^{4} z^{2} \left(37800 - 53550 a\right) + y^{4} \left(3132 a + 4638\right) + y^{3} z^{6} \left(5880 - 5880 a\right)\\
&\qquad + y^{3} z^{4} \left(183400 - 70000 a\right) + y^{3} z^{2} \left(214236 a + 205900\right) + y^{3} \left(20120 a - 52640\right)\\
&\qquad + 630 y^{2} z^{8} + y^{2} z^{6} \left(52920 a + 39200\right) + y^{2} z^{4} \left(1150100 a - 361850\right)\\
&\qquad + y^{2} z^{2} \left(1012212 a - 1608756\right) + y^{2} \left(142424 - 210592 a\right) + y z^{8} \left(6930 a - 3780\right)\\
&\qquad + y z^{6} \left(312620 a - 395480\right) + y z^{4} \left(- 10400 a - 2878100\right) + y z^{2} \left(874278 - 3116768 a\right)\\
&\qquad + y \left(287148 a + 150558\right) + z^{8} \left(- 1890 a - 13090\right) + z^{6} \left(- 476020 a - 114340\right)\\
&\qquad + z^{4} \left(2482481 - 3435465 a\right) + z^{2} \left(586308 a + 1916345\right) + 141719 a - 362989\Bigr]\\
&{}+x^{7}\Bigl[- 840 a y^{6} z + y^{5} z^{3} \left(7560 - 7560 a\right) + y^{5} z \left(21420 - 6300 a\right) + 8820 y^{4} z^{5}\\
&\qquad + y^{4} z^{3} \left(113400 a + 70000\right) + y^{4} z \left(210084 a - 107134\right) + 1680 a y^{3} z^{7}\\
&\qquad + y^{3} z^{5} \left(184240 a - 105840\right) + y^{3} z^{3} \left(1051000 a - 1534200\right) + y^{3} z \left(- 397888 a - 674268\right)\\
&\qquad + y^{2} z^{7} \left(12600 a - 27720\right) + y^{2} z^{5} \left(- 249780 a - 937860\right) + y^{2} z^{3} \left(27400 - 5784900 a\right)\\
&\qquad + y^{2} z \left(3112642 - 2236156 a\right) + y z^{7} \left(15400 - 120120 a\right) + y z^{5} \left(2867220 - 3546060 a\right)\\
&\qquad + y z^{3} \left(13746980 - 3789088 a\right) + y z \left(4991608 a - 1182968\right) - 70 z^{9} + z^{7} \left(199920 - 66220 a\right)\\
&\qquad + z^{5} \left(2066230 a + 2183680\right) + z^{3} \left(12151256 a - 6212842\right) + z \left(- 785906 a - 2703828\right)\Bigr]\\
&{}+x^{8}\Bigl[y^{7} \left(120 - 120 a\right) + 3780 y^{6} z^{2} + y^{6} \left(2520 a + 1050\right) + 8820 a y^{5} z^{4}\\
&\qquad + y^{5} z^{2} \left(110040 a - 68040\right) + y^{5} \left(20590 a - 42020\right) + y^{4} z^{6} \left(2940 a - 2940\right)\\
&\qquad + y^{4} z^{4} \left(98000 a - 230300\right) + y^{4} z^{2} \left(- 362950 a - 788250\right) + y^{4} \left(99520 - 267952 a\right)\\
&\qquad + y^{3} z^{6} \left(- 35280 a - 29400\right) + y^{3} z^{4} \left(418300 - 1981400 a\right)\\
&\qquad + y^{3} z^{2} \left(5792800 - 5758800 a\right) + y^{3} \left(292900 a + 743274\right) + y^{2} z^{6} \left(421400 - 366520 a\right)\\
&\qquad + y^{2} z^{4} \left(8874900 - 1679100 a\right) + y^{2} z^{2} \left(14978766 a + 5649384\right) + y^{2} \left(1892310 a - 2488988\right)\\
&\qquad - 560 a y z^{8} + y z^{6} \left(945980 a + 457660\right) + y z^{4} \left(21318610 a - 10445050\right)\\
&\qquad + y z^{2} \left(17722764 a - 36463092\right) + y \left(801004 - 3475076 a\right) + z^{8} \left(3290 - 1260 a\right)\\
&\qquad + z^{6} \left(1161050 a - 1627570\right) + z^{4} \left(- 3444640 a - 14775140\right) + z^{2} \left(8544658 - 25600678 a\right)\\
&\qquad + 513084 a + 1655387\Bigr]\\
&{}+x^{9}\Bigl[1080 a y^{7} z + y^{6} z^{3} \left(5880 a - 5880\right) + y^{6} z \left(14000 a - 36680\right) - 3528 y^{5} z^{5}\\
&\qquad + y^{5} z^{3} \left(- 105840 a - 78400\right) + y^{5} z \left(145400 - 460700 a\right) + y^{4} z^{5} \left(52920 - 97020 a\right)\\
&\qquad + y^{4} z^{3} \left(1983400 - 1563100 a\right) + y^{4} z \left(20300 a + 2880050\right) + y^{3} z^{5} \left(111720 a + 733040\right)\\
&\qquad + y^{3} z^{3} \left(9631400 a + 2214800\right) + y^{3} z \left(13757220 a - 10013780\right) + y^{2} z^{7} \left(1960 - 1960 a\right)\\
&\qquad + y^{2} z^{5} \left(4223520 a - 2868180\right) + y^{2} z^{3} \left(21657360 a - 42775340\right)\\
&\qquad + y^{2} z \left(- 18841866 a - 17630430\right) + y z^{7} \left(14560 a + 8680\right) + y z^{5} \left(- 2928540 a - 6935160\right)\\
&\qquad + y z^{3} \left(14365760 - 73365260 a\right) + y z \left(51189024 - 33903048 a\right) + z^{7} \left(41300 a - 59220\right)\\
&\qquad + z^{5} \left(7559050 - 8869506 a\right) + z^{3} \left(54487632 - 7161702 a\right) + z \left(29896586 a - 5286324\right)\Bigr]\\
&{}+x^{10}\Bigl[y^{8} \left(135 a - 135\right) - 2520 y^{7} z^{2} + y^{7} \left(- 3240 a - 2000\right) - 2940 a y^{6} z^{4}\\
&\qquad + y^{6} z^{2} \left(52920 - 92120 a\right) + y^{6} \left(76820 - 52550 a\right) + y^{5} z^{4} \left(97020 - 44100 a\right)\\
&\qquad + y^{5} z^{2} \left(253380 a + 937860\right) + y^{5} \left(580860 a - 4720\right) + y^{4} z^{4} \left(1058400 a - 142100\right)\\
&\qquad + y^{4} z^{2} \left(8894400 a - 7251300\right) + y^{4} \left(930160 a - 3440585\right) + 3920 y^{3} z^{6}\\
&\qquad + y^{3} z^{4} \left(2229500 a - 7060200\right) + y^{3} z^{2} \left(- 21345940 a - 21567520\right)\\
&\qquad + y^{3} \left(6314380 - 12160460 a\right) + y^{2} z^{6} \left(70560 a - 45080\right) + y^{2} z^{4} \left(7645290 - 24904950 a\right)\\
&\qquad + y^{2} z^{2} \left(110776650 - 88346940 a\right) + y^{2} \left(8740740 a + 16847244\right) + y z^{6} \left(- 119140 a - 248640\right)\\
&\qquad + y z^{4} \left(44462990 - 5745780 a\right) + y z^{2} \left(143679990 a + 19834860\right) + y \left(24411924 a - 29861264\right)\\
&\qquad + z^{8} \left(140 a - 140\right) + z^{6} \left(556220 - 576660 a\right) + z^{4} \left(38669030 a - 18846150\right)\\
&\qquad + z^{2} \left(44597758 a - 116164920\right) - 15005311 a - 544778\Bigr]\\
&{}+x^{11}\Bigl[- 630 a y^{8} z + y^{7} z^{3} \left(1680 - 1680 a\right) + y^{7} z \left(26320 - 11200 a\right)\\
&\qquad + y^{6} z^{3} \left(35280 a + 29400\right) + y^{6} z \left(397480 a - 84860\right) + y^{5} z^{3} \left(733040 a - 848680\right)\\
&\qquad + y^{5} z \left(650040 a - 3561660\right) + 4900 a y^{4} z^{5} + y^{4} z^{3} \left(- 4881100 a - 2200100\right)\\
&\qquad + y^{4} z \left(10786990 - 21525790 a\right) + y^{3} z^{5} \left(41160 a - 119560\right)\\
&\qquad + y^{3} z^{3} \left(33534200 - 22908040 a\right) + y^{3} z \left(15540700 a + 58795940\right)\\
&\qquad + y^{2} z^{5} \left(338100 - 962640 a\right) + y^{2} z^{3} \left(79280140 a + 9872080\right)\\
&\qquad + y^{2} z \left(163563900 a - 145378620\right) - 840 y z^{7} + y z^{5} \left(20720 a + 2994600\right)\\
&\qquad + y z^{3} \left(76662620 a - 156536960\right) + y z \left(- 146104300 a - 86754860\right) + z^{7} \left(3920 - 5460 a\right)\\
&\qquad + z^{5} \left(4513530 a - 2993810\right) + z^{3} \left(14467230 - 101150360 a\right) + z \left(135041266 - 79575830 a\right)\Bigr]\\
&{}+x^{12}\Bigl[y^{9} \left(70 - 70 a\right) + 630 y^{8} z^{2} + y^{8} \left(1890 a + 1400\right) + y^{7} z^{2} \left(27720 a - 15120\right)\\
&\qquad + y^{7} \left(44660 a - 56840\right) + y^{6} z^{2} \left(- 58800 a - 366520\right) + y^{6} \left(- 487120 a - 107140\right)\\
&\qquad + y^{5} z^{4} \left(3920 a - 3920\right) + y^{5} z^{2} \left(2958900 - 4261320 a\right) + y^{5} \left(4318970 - 2138760 a\right)\\
&\qquad + y^{4} z^{4} \left(- 83300 a - 39200\right) + y^{4} z^{2} \left(8294340 a + 17102010\right) + y^{4} \left(18705470 a - 4032110\right)\\
&\qquad + y^{3} z^{4} \left(1369200 - 838600 a\right) + y^{3} z^{2} \left(89373780 a - 81127660\right) + y^{3} \left(5512120 a - 54536680\right)\\
&\qquad - 2100 a y^{2} z^{6} + y^{2} z^{4} \left(6722800 a - 403900\right) + y^{2} z^{2} \left(- 126562290 a - 111012540\right)\\
&\qquad + y^{2} \left(74530190 - 116689560 a\right) + y z^{6} \left(28140 - 7140 a\right) + y z^{4} \left(5749280 a - 19792780\right)\\
&\qquad + y z^{2} \left(311292480 - 255832080 a\right) + y \left(57536360 a + 78070540\right) + z^{6} \left(87780 a - 44660\right)\\
&\qquad + z^{4} \left(9049550 - 21721090 a\right) + z^{2} \left(153817680 a + 43554560\right) + 51428936 a - 66560746\Bigr]\\
&{}+x^{13}\Bigl[140 a y^{9} z + y^{8} z \left(3150 a - 6930\right) + y^{7} z \left(17080 - 121800 a\right) - 1960 y^{6} z^{3}\\
&\qquad + y^{6} z \left(1424640 - 428260 a\right) + y^{5} z^{3} \left(54880 - 76440 a\right) + y^{5} z \left(10256820 a - 3447780\right)\\
&\qquad + y^{4} z^{3} \left(497000 a + 636300\right) + y^{4} z \left(3306840 a - 44795330\right) + y^{3} z^{5} \left(2800 - 2800 a\right)\\
&\qquad + y^{3} z^{3} \left(6862800 a - 7862400\right) + y^{3} z \left(88886160 - 160214920 a\right) + y^{2} z^{5} \left(46200 a + 12600\right)\\
&\qquad + y^{2} z^{3} \left(- 26146380 a - 7826860\right) + y^{2} z \left(67566960 a + 251519220\right) + y z^{5} \left(170520 a - 388920\right)\\
&\qquad + y z^{3} \left(77719460 - 40983200 a\right) + y z \left(392466420 a - 322381340\right) + 140 a z^{7}\\
&\qquad + z^{5} \left(258860 - 786240 a\right) + z^{3} \left(65853920 a - 11843680\right) + z \left(- 119811310 a - 116823990\right)\Bigr]\\
&{}+x^{14}\Bigl[y^{10} \left(14 a - 14\right) + y^{9} \left(- 420 a - 350\right) + y^{8} \left(15260 - 13090 a\right) - 560 a y^{7} z^{2}\\
&\qquad + y^{7} \left(143780 a + 60340\right) + y^{6} z^{2} \left(27440 - 5880 a\right) + y^{6} \left(1129520 a - 1725850\right)\\
&\qquad + y^{5} z^{2} \left(538020 a - 277620\right) + y^{5} \left(- 8482810 a - 497476\right) + 2100 y^{4} z^{4}\\
&\qquad + y^{4} z^{2} \left(- 1089200 a - 3971100\right) + y^{4} \left(40508020 - 17155960 a\right) + y^{3} z^{4} \left(63000 a - 53200\right)\\
&\qquad + y^{3} z^{2} \left(24116120 - 28298340 a\right) + y^{3} \left(108944120 a - 26573740\right)\\
&\qquad + y^{2} z^{4} \left(- 429660 a - 248640\right) + y^{2} z^{2} \left(55559100 a + 46926900\right)\\
&\qquad + y^{2} \left(26484020 a - 195006220\right) + y z^{6} \left(560 a - 560\right) + y z^{4} \left(2954840 - 1747340 a\right)\\
&\qquad + y z^{2} \left(135349780 a - 181368660\right) + y \left(130474810 - 237389530 a\right) + z^{6} \left(- 3780 a - 560\right)\\
&\qquad + z^{4} \left(4358195 a - 708675\right) + z^{2} \left(- 122780160 a - 9142480\right) + 32610896 a + 85823498\Bigr]\\
&{}+x^{15}\Bigl[y^{8} z \left(70 - 70 a\right) + y^{7} z \left(4480 a + 280\right) + y^{6} z \left(46340 a - 131880\right) + 840 a y^{5} z^{3}\\
&\qquad + y^{5} z \left(568960 - 1685320 a\right) + y^{4} z^{3} \left(2100 a - 35700\right) + y^{4} z \left(11280430 - 282380 a\right)\\
&\qquad + y^{3} z^{3} \left(428120 - 579880 a\right) + y^{3} z \left(58698080 a - 37107000\right) - 840 y^{2} z^{5}\\
&\qquad + y^{2} z^{3} \left(2153340 a + 2046520\right) + y^{2} z \left(- 55844860 a - 108489920\right) + y z^{5} \left(13440 - 14280 a\right)\\
&\qquad + y z^{3} \left(9787120 a - 13489000\right) + y z \left(232809100 - 225909200 a\right) + z^{5} \left(46410 a + 12460\right)\\
&\qquad + z^{3} \left(38260 - 15464760 a\right) + z \left(128584690 a + 47607820\right)\Bigr]\\
&{}+x^{16}\Bigl[y^{8} \left(210 a - 350\right) + y^{7} \left(- 11200 a - 840\right) + y^{6} z^{2} \left(140 a - 140\right)\\
&\qquad + y^{6} \left(223720 - 107100 a\right) + y^{5} z^{2} \left(1260 - 10920 a\right) + y^{5} \left(1978620 a - 302540\right)\\
&\qquad + y^{4} z^{2} \left(242340 - 21840 a\right) + y^{4} \left(2975840 a - 12210450\right) - 560 a y^{3} z^{4}\\
&\qquad + y^{3} z^{2} \left(2736860 a - 1806140\right) + y^{3} \left(21668920 - 48820380 a\right) + y^{2} z^{4} \left(840 a + 16800\right)\\
&\qquad + y^{2} z^{2} \left(- 6132000 a - 8696520\right) + y^{2} \left(15559380 a + 93710320\right) + y z^{4} \left(158130 a - 144130\right)\\
&\qquad + y z^{2} \left(37263660 - 31518060 a\right) + y \left(154112960 a - 126732800\right) + 70 z^{6}\\
&\qquad + z^{4} \left(- 333830 a - 118860\right) + z^{2} \left(34578030 a + 5182450\right) - 57839390 a + 42459540\Bigr]\\
&{}+x^{17}\Bigl[y^{6} z \left(1400 - 840 a\right) + y^{5} z \left(51240 a - 10920\right) + y^{4} z^{3} \left(140 - 140 a\right)\\
&\qquad + y^{4} z \left(74340 a - 762300\right) + y^{3} z^{3} \left(10080 a - 1960\right) + y^{3} z \left(4000640 - 6639360 a\right)\\
&\qquad + y^{2} z^{3} \left(- 16100 a - 141680\right) + y^{2} z \left(9412620 a + 18990720\right) + 140 a y z^{5}\\
&\qquad + y z^{3} \left(871640 - 959000 a\right) + y z \left(55149080 a - 57934520\right) + z^{5} \left(- 210 a - 1610\right)\\
&\qquad + z^{3} \left(1521520 a + 608440\right) + z \left(- 44786060 a - 12873480\right)\Bigr]\\
&{}+x^{18}\Bigl[y^{6} \left(700 a - 2940\right) + y^{5} \left(29820 - 86520 a\right) + y^{4} z^{2} \left(1260 a - 2100\right)\\
&\qquad + y^{4} \left(941640 - 79800 a\right) + y^{3} z^{2} \left(24360 - 68880 a\right) + y^{3} \left(6626340 a - 3712940\right)\\
&\qquad + y^{2} z^{4} \left(70 a - 70\right) + y^{2} z^{2} \left(128940 a + 612360\right) + y^{2} \left(- 6090000 a - 16976680\right)\\
&\qquad + y z^{4} \left(770 - 3220 a\right) + y z^{2} \left(3336060 a - 3107580\right) + y \left(39160240 - 40771780 a\right)\\
&\qquad + z^{4} \left(4690 a + 15960\right) + z^{2} \left(- 4383120 a - 1712200\right) + 25782295 a + 10347625\Bigr]\\
&{}+x^{19}\Bigl[y^{4} z \left(8820 - 2100 a\right) + y^{3} z \left(213920 a - 110040\right) + y^{2} z^{3} \left(1400 - 840 a\right)\\
&\qquad + y^{2} z \left(- 469980 a - 1331400\right) + y z^{3} \left(28840 a - 12600\right) + y z \left(6142360 - 6261360 a\right)\\
&\qquad + z^{5} \left(14 - 14 a\right) + z^{3} \left(- 45780 a - 83300\right) + z \left(7302330 a + 2428860\right)\Bigr]\\
&{}+x^{20}\Bigl[y^{4} \left(- 2380 a - 10220\right) + y^{3} \left(175000 - 248640 a\right) + y^{2} z^{2} \left(2100 a - 8820\right)\\
&\qquad + y^{2} \left(648620 a + 1147160\right) + y z^{2} \left(80220 - 127400 a\right) + y \left(4905810 a - 5199810\right)\\
&\qquad + z^{4} \left(210 a - 350\right) + z^{2} \left(228900 a + 226310\right) - 5332390 a + 1292060\Bigr]\\
&{}+x^{21}\Bigl[y^{2} z \left(4760 a + 20440\right) + y z \left(274540 a - 233240\right) + z^{3} \left(2940 - 700 a\right)\\
&\qquad + z \left(- 563290 a - 272370\right)\Bigr]\\
&{}+x^{22}\Bigl[y^{2} \left(- 16170 a - 12950\right) + y \left(255570 - 223300 a\right) + z^{2} \left(- 2380 a - 10220\right) + 536970 a + 74620\Bigr]\\
&{}+x^{23}\Bigl[z \left(16170 a + 12950\right)\Bigr]\\
&{}+x^{24}\Bigl[- 21350 a + 1302\Bigr]
\end{align*}
\endgroup

\subsubsection*{Inverse map}
\paragraph{Component 1.}
\begingroup
\allowdisplaybreaks[0]
\begin{gather*}
(\Phi_T^{-1})_1(X,Y,Z)=X^3-3X+Y.
\end{gather*}
\endgroup

\paragraph{Component 2.}
\begingroup
\footnotesize
\allowdisplaybreaks[0]
\begin{align*}
(\Phi_T^{-1})_2(X,Y,Z)={}&X^{6} \left(42 a - 42\right) + X^{4} \left(175 - 175 a\right) + X^{2} \left(147 a - 162\right) + 8\\
&{}+(X^3-3X+Y)\Bigl[- 48 X^{7} a + 280 X^{5} a + X^{3} \left(20 - 466 a\right) + X \left(165 a - 50\right)\Bigr]\\
&{}+(X^3-3X+Y)^{2}\Bigl[81 X^{8} - 490 X^{6} + X^{4} \left(45 a + 816\right) + X^{2} \left(- 125 a - 225\right) + 4 a - 6\Bigr]\\
&{}+(X^3-3X+Y)^{3}\Bigl[X^{9} \left(140 a - 140\right) + X^{7} \left(920 - 920 a\right) + X^{5} \left(1726 a - 1846\right)\\
&\qquad + X^{3} \left(1070 - 670 a\right) + X \left(- 49 a - 91\right)\Bigr]\\
&{}+(X^3-3X+Y)^{4}\Bigl[- 196 X^{10} a + 1400 X^{8} a + X^{6} \left(140 - 3312 a\right) + X^{4} \left(2900 a - 500\right)\\
&\qquad + X^{2} \left(265 - 816 a\right)\Bigr]\\
&{}+(Z+8-14X^2)\Bigl[- 1\Bigr]\\
&{}+(Z+8-14X^2)(X^3-3X+Y)\Bigl[X \left(2 a - 2\right)\Bigr]\\
&{}+(Z+8-14X^2)(X^3-3X+Y)^{2}\Bigl[- 2 X^{2} a + 5 a\Bigr]\\
&{}+(Z+8-14X^2)(X^3-3X+Y)^{3}\Bigl[4 X^{3} - 10 X\Bigr]\\
&{}+(Z+8-14X^2)(X^3-3X+Y)^{4}\Bigl[X^{4} \left(10 a - 10\right) + X^{2} \left(30 - 30 a\right) - a + 4\Bigr]\\
&{}+(Z+8-14X^2)(X^3-3X+Y)^{5}\Bigl[- 28 X^{5} a + 100 X^{3} a + X \left(10 - 58 a\right)\Bigr]\\
&{}+(Z+8-14X^2)^{2}(X^3-3X+Y)^{6}\Bigl[- a\Bigr]
\end{align*}
\endgroup

\paragraph{Component 3.}
\begingroup
\footnotesize
\allowdisplaybreaks[0]
\begin{align*}
(\Phi_T^{-1})_3(X,Y,Z)={}&- X\\
&{}+(X^3-3X+Y)\Bigl[42 X^{6} - 175 X^{4} + X^{2} \left(14 a + 148\right) - 8 a - 3\Bigr]\\
&{}+(X^3-3X+Y)^{2}\Bigl[X^{7} \left(48 a - 48\right) + X^{5} \left(280 - 280 a\right) + X^{3} \left(448 a - 466\right)\\
&\qquad + X \left(165 - 120 a\right)\Bigr]\\
&{}+(X^3-3X+Y)^{3}\Bigl[- 81 X^{8} a + 490 X^{6} a + X^{4} \left(40 - 861 a\right) + X^{2} \left(350 a - 110\right) + 4\Bigr]\\
&{}+(X^3-3X+Y)^{4}\Bigl[140 X^{9} - 920 X^{7} + X^{5} \left(106 a + 1740\right) + X^{3} \left(- 350 a - 720\right)\\
&\qquad + X \left(116 a - 20\right)\Bigr]\\
&{}+(X^3-3X+Y)^{5}\Bigl[X^{10} \left(196 a - 196\right) + X^{8} \left(1400 - 1400 a\right) + X^{6} \left(3172 a - 3312\right)\\
&\qquad + X^{4} \left(2900 - 2400 a\right) + X^{2} \left(551 a - 816\right)\Bigr]\\
&{}+(Z+8-14X^2)(X^3-3X+Y)\Bigl[a\Bigr]\\
&{}+(Z+8-14X^2)(X^3-3X+Y)^{2}\Bigl[2 X\Bigr]\\
&{}+(Z+8-14X^2)(X^3-3X+Y)^{3}\Bigl[X^{2} \left(2 a - 2\right) - 5 a - 5\Bigr]\\
&{}+(Z+8-14X^2)(X^3-3X+Y)^{4}\Bigl[- 4 X^{3} a + 10 X a\Bigr]\\
&{}+(Z+8-14X^2)(X^3-3X+Y)^{5}\Bigl[10 X^{4} - 30 X^{2} + 4 a\Bigr]\\
&{}+(Z+8-14X^2)(X^3-3X+Y)^{6}\Bigl[X^{5} \left(28 a - 28\right) + X^{3} \left(100 - 100 a\right) + X \left(48 a - 58\right)\Bigr]\\
&{}+(Z+8-14X^2)^{2}(X^3-3X+Y)^{7}\Bigl[a - 1\Bigr]
\end{align*}
\endgroup



\clearpage
\section{The harmonic figure-eight}

The second curve is
\begin{equation}\label{eq:harmonic}
\gamma_H(t)=
(t^3-3t,\ t^5-5t^3+5t,\ t^7-7t^5+14t^3-7t).
\end{equation}
The coefficient in the second component is $+5t$.
Indeed,
\[
\gamma_H(t)=\bigl(2T_3(t/2),2T_5(t/2),2T_7(t/2)\bigr),
\]
where $T_n(\cos\theta)=\cos(n\theta)$. This identifies the parametrization
with the harmonic knot $H(3,5,7)$ of \cite{KP} by linear scalings.
It is distinct from the figure-eight parametrization in
\cite[\S3, Example~2]{Shastri}, whose second component contains $+4t$.

\subsection{Modification and correction}

Take $m(x)=x$ and $p(x,y)=2-y^2$. The modification
\[
\beta(x,y,z)=(x,xz-y^2+2,y)
\]
satisfies
\begin{equation}\label{eq:curve-transition}
\boxed{\beta(\gamma_H(t))=M(\gamma_T(t)).}
\end{equation}
Its inverse after adjoining $1/x$ is
\[
\beta^{-1}(x,Y,Z)=\left(x,Z,\frac{Y+Z^2-2}{x}\right).
\]
The hypothesis of Theorem~\ref{thm:transfer} follows from
\[
H(0,y,z)=-z,\qquad H(0,2-y^2,y)=-y.
\]
Thus the centre inverse is $\psi(T)=-T$, and the correction is
\begin{equation}\label{eq:harmonic-correction}
\begin{aligned}
\varphi(T)&=W(0,2-T^2,-T)\\
&=42(a-1)T^6+175(1-a)T^4+(147a-161)T^2+8.
\end{aligned}
\end{equation}
The transferred coordinates are
\begin{equation}\label{eq:harmonic-relative}
\begin{aligned}
h_H(x,y,z)&=H(x,xz-y^2+2,y),\\
q_H(x,y,z)&=
\frac{W(x,xz-y^2+2,y)-\varphi(h_H(x,y,z))}{x}.
\end{aligned}
\end{equation}
The inverse of $(x,h_H,q_H)$ has components
\begin{equation}\label{eq:harmonic-relative-inverse}
\begin{aligned}
R_H(x,T,Q)&=S(x,T,xQ+\varphi(T)),\\
S_H(x,T,Q)&=
\frac{\begin{gathered}
R(x,T,xQ+\varphi(T))\\[-2pt]
+S(x,T,xQ+\varphi(T))^2-2
\end{gathered}}{x}.
\end{aligned}
\end{equation}
Both numerators divided by $x$ are exactly divisible. This is established
by Theorem~\ref{thm:transfer} and verified directly in the certificate.

\subsection{The transition is global}

Over $K[x,1/x]$, set
\[
C(x,T,V)=\left(x,T,\frac{V-\varphi(T)}x\right),\qquad
C^{-1}(x,T,Q)=(x,T,xQ+\varphi(T)).
\]
Then
\begin{equation}\label{eq:harmonic-factorization}
\begin{aligned}
(x,h_H,q_H)&=C\comp A_0\comp\beta,\\
(x,R_H,S_H)&=\beta^{-1}\comp A_0^{-1}\comp C^{-1}.
\end{aligned}
\end{equation}
The individual maps $\beta^{-1}$ and $C$ require $1/x$. The complete maps
in \eqref{eq:harmonic-factorization} are polynomial by the two exact
divisions. Injectivity of localization therefore proves both composition
identities in $K[x,y,z]$, including $x=0$.

\subsection{Final rectifier and inverse}

Curve substitution gives
\[
h_H(\gamma_H(t))=t,\qquad
q_H(\gamma_H(t))=(1-a)(42t^3-49t).
\]
Define
\begin{equation}\label{eq:PhiH}
\begin{aligned}
(\Phi_H)_1(x,y,z)&=h_H(x,y,z),\\
(\Phi_H)_2(x,y,z)&=x-h_H(x,y,z)^3+3h_H(x,y,z),\\
(\Phi_H)_3(x,y,z)&=q_H(x,y,z)\\
&\quad-(1-a)\bigl(42h_H(x,y,z)^3-49h_H(x,y,z)\bigr).
\end{aligned}
\end{equation}
The inverse is
\begin{equation}\label{eq:PhiHinv}
\begin{aligned}
(\Phi_H^{-1})_1(X,Y,Z)&=X^3-3X+Y,\\
(\Phi_H^{-1})_2(X,Y,Z)&=
R_H\bigl(X^3-3X+Y,X,Z+(1-a)(42X^3-49X)\bigr),\\
(\Phi_H^{-1})_3(X,Y,Z)&=
S_H\bigl(X^3-3X+Y,X,Z+(1-a)(42X^3-49X)\bigr).
\end{aligned}
\end{equation}
All auxiliary polynomials in these formulas have been explicitly defined
in \eqref{eq:H}--\eqref{eq:RS} and
\eqref{eq:harmonic-relative}--\eqref{eq:harmonic-relative-inverse}.
The final transition is the permutation and shear
\[
L_H(x,T,Q)=(T,x-T^3+3T,Q-(1-a)(42T^3-49T)),
\]
so
\[
\boxed{\Phi_H=L_H\comp C\comp A_0\comp\beta.}
\]
The resulting identities and degrees are
\begin{equation}\label{eq:harmonic-result}
\boxed{\Phi_H(\gamma_H(t))=(t,0,0),\qquad
\deg\Phi_H=(7,21,41),\qquad \det D\Phi_H=-1.}
\end{equation}

\clearpage
\subsection{Coordinate-only formulas}

The following formulas use only the input coordinates and the constant
$a=(1+i\sqrt3)/2$. Repeated expressions are written in full.

\subsubsection*{Forward map}

\paragraph{Component 1.}

\begingroup\small\allowdisplaybreaks[4]

\begin{align*}
(\Phi_H)_1(x,y,z)={}&-y+x\Bigl[(1-a)\bigl(ax(xz-y^2+2)+y\bigr)^2\\
&\qquad +\bigl((a-1)x^2-a\bigr)(xz-y^2+2)\\
&\qquad +ax(1-x^2)y-2x^4+(6-2a)x^2-2(1-a)\Bigr]
\end{align*}

\endgroup

\paragraph{Component 2.}

\begingroup\small\allowdisplaybreaks[4]

\begin{align*}
(\Phi_H)_2(x,y,z)={}&x-\Bigl\{-y+x\Bigl[(1-a)\bigl(ax(xz-y^2+2)+y\bigr)^2\\
&\qquad +\bigl((a-1)x^2-a\bigr)(xz-y^2+2)\\
&\qquad +ax(1-x^2)y-2x^4+(6-2a)x^2-2(1-a)\Bigr]\Bigr\}^{3}\\
&{}+3\Bigl\{-y+x\Bigl[(1-a)\bigl(ax(xz-y^2+2)+y\bigr)^2\\
&\qquad +\bigl((a-1)x^2-a\bigr)(xz-y^2+2)\\
&\qquad +ax(1-x^2)y-2x^4+(6-2a)x^2-2(1-a)\Bigr]\Bigr\}
\end{align*}

\endgroup

\clearpage

\paragraph{Component 3.}

\begingroup\footnotesize\allowdisplaybreaks[4]

\begin{align*}
(\Phi_H)_3(x,y,z)={}&y^{3} \left(42 - 42 a\right) + y \left(49 a - 49\right) - z\\
&{}- 120 a y^{7} + 574 a y^{5} + y^{3} \left(- 934 a - 20\right) + y \left(463 a + 10\right)\\
&{}+(xz-y^2+2)\Bigl[- 252 y^{5} + 700 y^{3} + y \left(- 28 a - 296\right)\Bigr]\\
&{}+x\Bigl[135 y^{8} - 630 y^{6} + y^{4} \left(40 a + 1112\right) + y^{2} \left(- 20 a - 1129\right)\\
&\qquad - 26 a - 352\Bigr]\\
&{}+x(xz-y^2+2)\Bigl[y^{6} \left(840 - 840 a\right) + y^{4} \left(2870 a - 2870\right)\\
&\qquad + y^{2} \left(2802 - 2858 a\right) + 472 a - 463\Bigr]\\
&{}+x(xz-y^2+2)^{2}\Bigl[- 630 a y^{4} + 1050 a y^{2} - 162 a - 13\Bigr]\\
&{}+x^{2}\Bigl[y^{9} \left(70 a - 70\right) + y^{7} \left(160 - 160 a\right)\\
&\qquad + y^{5} \left(830 - 1254 a\right) + y^{3} \left(3784 a - 2444\right)\\
&\qquad + y \left(711 - 1331 a\right)\Bigr]\\
&{}+x^{2}(xz-y^2+2)\Bigl[1080 a y^{7} - 3780 a y^{5} + y^{3} \left(4608 a - 152\right)\\
&\qquad + y \left(36 - 2298 a\right)\Bigr]\\
&{}+x^{2}(xz-y^2+2)^{2}\Bigl[2520 y^{5} - 5740 y^{3} + y \left(2854 - 52 a\right)\Bigr]\\
&{}+x^{2}(xz-y^2+2)^{3}\Bigl[y^{3} \left(840 - 840 a\right) + y \left(700 a - 700\right)\Bigr]\\
&{}+x^{3}\Bigl[- 14 a y^{10} - 140 a y^{8} + y^{6} \left(2330 a + 580\right)\\
&\qquad + y^{4} \left(- 7550 a - 2400\right) + y^{2} \left(9147 a + 4135\right) - 1346 a + 940\Bigr]\\
&{}+x^{3}(xz-y^2+2)\Bigl[- 630 y^{8} + 1120 y^{6} + y^{4} \left(6286 - 2136 a\right)\\
&\qquad + y^{2} \left(4032 a - 11364\right) - 620 a - 1332\Bigr]\\
&{}+x^{3}(xz-y^2+2)^{2}\Bigl[y^{6} \left(3780 a - 3780\right) + y^{4} \left(9450 - 9450 a\right)\\
&\qquad + y^{2} \left(6696 a - 6912\right) - 1133 a - 1149\Bigr]\\
&{}+x^{3}(xz-y^2+2)^{3}\Bigl[4200 a y^{4} - 5740 a y^{2} + 934 a + 16\Bigr]\\
&{}+x^{3}(xz-y^2+2)^{4}\Bigl[630 y^{2} - 175\Bigr]\\
&{}+x^{4}\Bigl[70 y^{9} + y^{7} \left(160 a - 1420\right) + y^{5} \left(190 a + 3530\right)\\
&\qquad + y^{3} \left(- 2096 a - 3856\right) + y \left(10532 - 2220 a\right)\Bigr]\\
&{}+x^{4}(xz-y^2+2)\Bigl[y^{9} \left(140 - 140 a\right) + y^{7} \left(1120 - 1120 a\right)\\
&\qquad + y^{5} \left(17680 a - 13980\right) + y^{3} \left(30200 - 40468 a\right)\\
&\qquad + y \left(26882 a - 18320\right)\Bigr]\\
&{}+x^{4}(xz-y^2+2)^{2}\Bigl[- 2520 a y^{7} + 3360 a y^{5} + y^{3} \left(8300 a + 4304\right)\\
&\qquad + y \left(- 7332 a - 4044\right)\Bigr]\\
&{}+x^{4}(xz-y^2+2)^{3}\Bigl[- 7560 y^{5} + 12600 y^{3} + y \left(- 136 a - 4472\right)\Bigr]\\
&{}+x^{4}(xz-y^2+2)^{4}\Bigl[y^{3} \left(4200 a - 4200\right) + y \left(2870 - 2870 a\right)\Bigr]\\
&{}+x^{4}(xz-y^2+2)^{5}\Bigl[252 a y\Bigr]\\
&{}+x^{5}\Bigl[140 y^{8} + y^{6} \left(1680 - 3940 a\right) + y^{4} \left(25535 a - 11419\right)\\
&\qquad + y^{2} \left(16523 - 43464 a\right) - 336 a - 2356\Bigr]\\
&{}+x^{5}(xz-y^2+2)\Bigl[630 a y^{8} + y^{6} \left(- 8820 a - 1520\right) + y^{4} \left(18600 a + 600\right)\\
&\qquad + y^{2} \left(4002 - 17748 a\right) + 8304 a + 2692\Bigr]\\
&{}+x^{5}(xz-y^2+2)^{2}\Bigl[630 y^{8} + 3920 y^{6} + y^{4} \left(9800 a - 44750\right)\\
&\qquad + y^{2} \left(61704 - 16404 a\right) + 4428 a - 13600\Bigr]\\
&{}+x^{5}(xz-y^2+2)^{3}\Bigl[y^{6} \left(5880 - 5880 a\right) + y^{4} \left(5600 a - 5600\right)\\
&\qquad + y^{2} \left(12636 a - 8300\right) - 3796 a - 2444\Bigr]\\
&{}+x^{5}(xz-y^2+2)^{4}\Bigl[- 9450 a y^{4} + 9450 a y^{2} - 1152 a - 32\Bigr]\\
&{}+x^{5}(xz-y^2+2)^{5}\Bigl[- 2520 y^{2} + 574\Bigr]\\
&{}+x^{5}(xz-y^2+2)^{6}\Bigl[42 a - 42\Bigr]\\
&{}+x^{6}\Bigl[- 70 y^{9} + y^{7} \left(980 a + 840\right) + y^{5} \left(- 9390 a - 5960\right)\\
&\qquad + y^{3} \left(26352 a + 18598\right) + y \left(- 29700 a - 42708\right)\Bigr]\\
&{}+x^{6}(xz-y^2+2)\Bigl[y^{7} \left(840 a + 280\right) + y^{5} \left(22740 - 12660 a\right)\\
&\qquad + y^{3} \left(55912 a - 101220\right) + y \left(85240 - 52272 a\right)\Bigr]\\
&{}+x^{6}(xz-y^2+2)^{2}\Bigl[y^{7} \left(2520 a - 2520\right) + y^{5} \left(26460 - 32220 a\right)\\
&\qquad + y^{3} \left(41500 a - 37200\right) + y \left(17646 - 15924 a\right)\Bigr]\\
&{}+x^{6}(xz-y^2+2)^{3}\Bigl[1680 a y^{7} + 7840 a y^{5} + y^{3} \left(- 46600 a - 13800\right)\\
&\qquad + y \left(30200 a + 11604\right)\Bigr]\\
&{}+x^{6}(xz-y^2+2)^{4}\Bigl[8820 y^{5} - 5600 y^{3} + y \left(2184 a - 6334\right)\Bigr]\\
&{}+x^{6}(xz-y^2+2)^{5}\Bigl[y^{3} \left(7560 - 7560 a\right) + y \left(3780 a - 3780\right)\Bigr]\\
&{}+x^{6}(xz-y^2+2)^{6}\Bigl[- 840 a y\Bigr]\\
&{}+x^{7}\Bigl[y^{8} \left(140 - 140 a\right) + y^{6} \left(4090 a - 2730\right)\\
&\qquad + y^{4} \left(29360 - 31050 a\right) + y^{2} \left(69544 a - 62488\right) + 19252 a - 33824\Bigr]\\
&{}+x^{7}(xz-y^2+2)\Bigl[- 560 a y^{8} + y^{6} \left(13020 a - 6860\right) + y^{4} \left(45550 - 81190 a\right)\\
&\qquad + y^{2} \left(143380 a - 75288\right) - 72656 a - 29088\Bigr]\\
&{}+x^{7}(xz-y^2+2)^{2}\Bigl[y^{6} \left(3920 a - 1960\right) + y^{4} \left(25200 a + 29400\right)\\
&\qquad + y^{2} \left(- 67434 a - 83016\right) + 16448 a + 25328\Bigr]\\
&{}+x^{7}(xz-y^2+2)^{3}\Bigl[- 5880 y^{6} + y^{4} \left(55700 - 11600 a\right) + y^{2} \left(7400 a - 44600\right)\\
&\qquad - 184 a - 6034\Bigr]\\
&{}+x^{7}(xz-y^2+2)^{4}\Bigl[y^{6} \left(2940 a - 2940\right) + y^{4} \left(9800 a - 9800\right)\\
&\qquad + y^{2} \left(34950 - 45850 a\right) + 10618 a - 7550\Bigr]\\
&{}+x^{7}(xz-y^2+2)^{5}\Bigl[8820 a y^{4} - 3360 a y^{2} - 830 a + 440\Bigr]\\
&{}+x^{7}(xz-y^2+2)^{6}\Bigl[3780 y^{2} - 630\Bigr]\\
&{}+x^{7}(xz-y^2+2)^{7}\Bigl[- 120 a - 120\Bigr]\\
&{}+x^{8}\Bigl[y^{7} \left(- 1260 a - 140\right) + y^{5} \left(18774 a - 3454\right)\\
&\qquad + y^{3} \left(17932 - 97842 a\right) + y \left(182792 a - 5992\right)\Bigr]\\
&{}+x^{8}(xz-y^2+2)\Bigl[840 y^{7} + y^{5} \left(7260 a - 19440\right) + y^{3} \left(108880 - 3660 a\right)\\
&\qquad + y \left(6696 a - 115536\right)\Bigr]\\
&{}+x^{8}(xz-y^2+2)^{2}\Bigl[y^{7} \left(1960 - 1960 a\right) + y^{5} \left(19320 a - 39900\right)\\
&\qquad + y^{3} \left(171260 - 83040 a\right) + y \left(80414 a - 151910\right)\Bigr]\\
&{}+x^{8}(xz-y^2+2)^{3}\Bigl[y^{5} \left(5880 a - 7840\right) + y^{3} \left(69800 a - 33600\right)\\
&\qquad + y \left(44620 - 99380 a\right)\Bigr]\\
&{}+x^{8}(xz-y^2+2)^{4}\Bigl[- 8820 a y^{5} + y^{3} \left(44100 a + 13600\right) + y \left(- 18600 a - 5250\right)\Bigr]\\
&{}+x^{8}(xz-y^2+2)^{5}\Bigl[- 3528 y^{5} - 7840 y^{3} + y \left(18560 - 4580 a\right)\Bigr]\\
&{}+x^{8}(xz-y^2+2)^{6}\Bigl[y^{3} \left(5880 a - 5880\right) + y \left(1120 - 1120 a\right)\Bigr]\\
&{}+x^{8}(xz-y^2+2)^{7}\Bigl[1080 a y\Bigr]\\
&{}+x^{9}\Bigl[y^{8} \left(140 a - 140\right) + y^{6} \left(3920 - 2240 a\right)\\
&\qquad + y^{4} \left(22830 a - 37710\right) + y^{2} \left(91720 - 61032 a\right) - 3984 a - 109376\Bigr]\\
&{}+x^{9}(xz-y^2+2)\Bigl[y^{6} \left(8120 - 9380 a\right) + y^{4} \left(82200 a - 95270\right)\\
&\qquad + y^{2} \left(299520 - 260790 a\right) + 164944 a - 174240\Bigr]\\
&{}+x^{9}(xz-y^2+2)^{2}\Bigl[1960 a y^{6} + y^{4} \left(- 19950 a - 15510\right) + y^{2} \left(139500 a - 90\right)\\
&\qquad - 46340 a - 504\Bigr]\\
&{}+x^{9}(xz-y^2+2)^{3}\Bigl[3920 y^{6} + y^{4} \left(- 34300 a - 33600\right) + y^{2} \left(85260 a + 94880\right)\\
&\qquad - 22560 a + 30920\Bigr]\\
&{}+x^{9}(xz-y^2+2)^{4}\Bigl[- 9800 y^{4} + y^{2} \left(24900 a - 50100\right) - 13540 a - 24615\Bigr]\\
&{}+x^{9}(xz-y^2+2)^{5}\Bigl[y^{4} \left(8820 - 8820 a\right) + y^{2} \left(35820 a - 26460\right) - 5080 a - 3720\Bigr]\\
&{}+x^{9}(xz-y^2+2)^{6}\Bigl[- 2940 a y^{4} - 3920 a y^{2} + 2330 a + 800\Bigr]\\
&{}+x^{9}(xz-y^2+2)^{7}\Bigl[- 2520 y^{2} + 160\Bigr]\\
&{}+x^{9}(xz-y^2+2)^{8}\Bigl[135 a - 135\Bigr]\\
&{}+x^{10}\Bigl[140 a y^{7} + y^{5} \left(9470 - 15050 a\right) + y^{3} \left(127960 a - 87110\right)\\
&\qquad + y \left(217496 - 362280 a\right)\Bigr]\\
&{}+x^{10}(xz-y^2+2)\Bigl[- 840 y^{7} + y^{5} \left(7000 a + 10080\right) + y^{3} \left(- 38420 a - 68640\right)\\
&\qquad + y \left(27440 a + 70600\right)\Bigr]\\
&{}+x^{10}(xz-y^2+2)^{2}\Bigl[y^{5} \left(26460 - 4200 a\right) + y^{3} \left(- 19060 a - 174280\right)\\
&\qquad + y \left(25680 a + 279960\right)\Bigr]\\
&{}+x^{10}(xz-y^2+2)^{3}\Bigl[y^{5} \left(1960 a - 1960\right) + y^{3} \left(12600 - 29240 a\right)\\
&\qquad + y \left(76380 a - 80780\right)\Bigr]\\
&{}+x^{10}(xz-y^2+2)^{4}\Bigl[4900 a y^{5} + y^{3} \left(34300 - 69300 a\right) + y \left(94510 a - 41110\right)\Bigr]\\
&{}+x^{10}(xz-y^2+2)^{5}\Bigl[y^{3} \left(- 7840 a - 1960\right) + y \left(- 10080 a - 9060\right)\Bigr]\\
&{}+x^{10}(xz-y^2+2)^{6}\Bigl[5880 y^{3} + y \left(3520 a - 12340\right)\Bigr]\\
&{}+x^{10}(xz-y^2+2)^{7}\Bigl[y^{3} \left(1680 - 1680 a\right) + y \left(1120 - 1120 a\right)\Bigr]\\
&{}+x^{10}(xz-y^2+2)^{8}\Bigl[- 630 a y\Bigr]\\
&{}+x^{11}\Bigl[y^{6} \left(560 a - 2940\right) + y^{4} \left(30900 - 1640 a\right)\\
&\qquad + y^{2} \left(27660 a - 96460\right) - 128224 a + 112576\Bigr]\\
&{}+x^{11}(xz-y^2+2)\Bigl[y^{6} \left(1260 a - 1260\right) + y^{4} \left(76020 - 33420 a\right)\\
&\qquad + y^{2} \left(142320 a - 390240\right) - 104520 a - 301480\Bigr]\\
&{}+x^{11}(xz-y^2+2)^{2}\Bigl[- 2100 a y^{6} + y^{4} \left(26600 a - 9800\right) + y^{2} \left(26760 - 93450 a\right)\\
&\qquad + 7320 a + 3200\Bigr]\\
&{}+x^{11}(xz-y^2+2)^{3}\Bigl[y^{4} \left(33600 a + 7000\right) + y^{2} \left(13300 - 196140 a\right) + 103960 a - 4880\Bigr]\\
&{}+x^{11}(xz-y^2+2)^{4}\Bigl[y^{2} \left(8010 - 9060 a\right) - 2340 a + 14800\Bigr]\\
&{}+x^{11}(xz-y^2+2)^{5}\Bigl[y^{4} \left(3920 a - 3920\right) + y^{2} \left(42420 - 21840 a\right) + 11880 a - 19790\Bigr]\\
&{}+x^{11}(xz-y^2+2)^{6}\Bigl[y^{2} \left(3920 - 5880 a\right) - 3040 a - 1680\Bigr]\\
&{}+x^{11}(xz-y^2+2)^{7}\Bigl[2520 a y^{2} - 1260 a + 560\Bigr]\\
&{}+x^{11}(xz-y^2+2)^{8}\Bigl[630 y^{2} + 140\Bigr]\\
&{}+x^{11}(xz-y^2+2)^{9}\Bigl[- 70 a - 70\Bigr]\\
&{}+x^{12}\Bigl[140 a y^{7} + y^{5} \left(3220 a - 3500\right) + y^{3} \left(96360 - 77320 a\right)\\
&\qquad + y \left(339080 a - 402240\right)\Bigr]\\
&{}+x^{12}(xz-y^2+2)\Bigl[y^{5} \left(- 10080 a - 2520\right) + y^{3} \left(68320 a + 8620\right)\\
&\qquad + y \left(7440 a - 55760\right)\Bigr]\\
&{}+x^{12}(xz-y^2+2)^{2}\Bigl[- 4200 y^{5} + y^{3} \left(77020 a + 71940\right) + y \left(- 238080 a - 149880\right)\Bigr]\\
&{}+x^{12}(xz-y^2+2)^{3}\Bigl[y^{5} \left(2800 - 2800 a\right) + y^{3} \left(11200 a - 14000\right)\\
&\qquad + y \left(16560 - 18800 a\right)\Bigr]\\
&{}+x^{12}(xz-y^2+2)^{4}\Bigl[y^{3} \left(36400 a - 30100\right) + y \left(99470 - 95040 a\right)\Bigr]\\
&{}+x^{12}(xz-y^2+2)^{5}\Bigl[1960 a y^{3} + y \left(2100 - 4620 a\right)\Bigr]\\
&{}+x^{12}(xz-y^2+2)^{6}\Bigl[- 1960 y^{3} + y \left(6860 a + 7560\right)\Bigr]\\
&{}+x^{12}(xz-y^2+2)^{7}\Bigl[y \left(1960 - 840 a\right)\Bigr]\\
&{}+x^{12}(xz-y^2+2)^{8}\Bigl[y \left(630 a - 630\right)\Bigr]\\
&{}+x^{12}(xz-y^2+2)^{9}\Bigl[140 a y\Bigr]\\
&{}+x^{13}\Bigl[560 y^{6} + y^{4} \left(- 11765 a - 13715\right) + y^{2} \left(31120 a + 61280\right)\\
&\qquad + 265312 a - 20976\Bigr]\\
&{}+x^{13}(xz-y^2+2)\Bigl[y^{6} \left(560 a - 560\right) + y^{4} \left(2800 a - 21700\right)\\
&\qquad + y^{2} \left(32680 a + 240600\right) - 74400 a + 214400\Bigr]\\
&{}+x^{13}(xz-y^2+2)^{2}\Bigl[y^{4} \left(15960 - 20160 a\right) + y^{2} \left(47160 a - 31920\right) + 37160 a - 69160\Bigr]\\
&{}+x^{13}(xz-y^2+2)^{3}\Bigl[- 7000 a y^{4} + y^{2} \left(141260 a - 70280\right) - 125680 a - 77280\Bigr]\\
&{}+x^{13}(xz-y^2+2)^{4}\Bigl[2100 y^{4} + 5600 a y^{2} - 5460 a + 2080\Bigr]\\
&{}+x^{13}(xz-y^2+2)^{5}\Bigl[y^{2} \left(2940 a - 18900\right) + 574 a + 19600\Bigr]\\
&{}+x^{13}(xz-y^2+2)^{6}\Bigl[y^{2} \left(1960 a - 1960\right) - 1400 a - 1470\Bigr]\\
&{}+x^{13}(xz-y^2+2)^{7}\Bigl[- 560 a y^{2} + 2100 a - 980\Bigr]\\
&{}+x^{13}(xz-y^2+2)^{8}\Bigl[140 a + 140\Bigr]\\
&{}+x^{13}(xz-y^2+2)^{9}\Bigl[- 70\Bigr]\\
&{}+x^{13}(xz-y^2+2)^{10}\Bigl[14 a - 14\Bigr]\\
&{}+x^{14}\Bigl[y^{5} \left(1470 a - 980\right) + y^{3} \left(17320 a - 43120\right)\\
&\qquad + y \left(356160 - 140480 a\right)\Bigr]\\
&{}+x^{14}(xz-y^2+2)\Bigl[2520 a y^{5} + y^{3} \left(25760 - 58800 a\right) + y \left(47040 - 18000 a\right)\Bigr]\\
&{}+x^{14}(xz-y^2+2)^{2}\Bigl[- 840 y^{5} + y^{3} \left(- 37380 a - 9800\right) + y \left(245840 a - 25520\right)\Bigr]\\
&{}+x^{14}(xz-y^2+2)^{3}\Bigl[y^{3} \left(13720 - 3080 a\right) + y \left(11760 a + 12320\right)\Bigr]\\
&{}+x^{14}(xz-y^2+2)^{4}\Bigl[y^{3} \left(6300 - 6300 a\right) + y \left(31920 a - 64470\right)\Bigr]\\
&{}+x^{14}(xz-y^2+2)^{5}\Bigl[840 a y^{3} + y \left(8680 a - 5320\right)\Bigr]\\
&{}+x^{14}(xz-y^2+2)^{6}\Bigl[y \left(- 4620 a - 560\right)\Bigr]\\
&{}+x^{14}(xz-y^2+2)^{7}\Bigl[- 840 y\Bigr]\\
&{}+x^{14}(xz-y^2+2)^{8}\Bigl[y \left(70 - 70 a\right)\Bigr]\\
&{}+x^{15}\Bigl[70 y^{6} + y^{4} \left(7000 a + 2100\right) + y^{2} \left(- 62720 a - 5720\right)\\
&\qquad - 235840 a - 127360\Bigr]\\
&{}+x^{15}(xz-y^2+2)\Bigl[y^{4} \left(1050 a - 1330\right) + y^{2} \left(- 63840 a - 72240\right)\\
&\qquad + 132160 a + 47200\Bigr]\\
&{}+x^{15}(xz-y^2+2)^{2}\Bigl[y^{4} \left(4200 a - 4200\right) + y^{2} \left(31080 - 19320 a\right) - 11360 a - 79680\Bigr]\\
&{}+x^{15}(xz-y^2+2)^{3}\Bigl[- 560 a y^{4} + y^{2} \left(33740 - 43820 a\right) + 57680 a - 68880\Bigr]\\
&{}+x^{15}(xz-y^2+2)^{4}\Bigl[y^{2} \left(1260 a + 840\right) + 16940 a - 3500\Bigr]\\
&{}+x^{15}(xz-y^2+2)^{5}\Bigl[2940 y^{2} - 6160 a + 5040\Bigr]\\
&{}+x^{15}(xz-y^2+2)^{6}\Bigl[y^{2} \left(140 a - 140\right) + 1120 a - 2520\Bigr]\\
&{}+x^{15}(xz-y^2+2)^{7}\Bigl[- 560 a - 560\Bigr]\\
&{}+x^{15}(xz-y^2+2)^{8}\Bigl[- 140 a\Bigr]\\
&{}+x^{16}\Bigl[y^{5} \left(490 - 490 a\right) + y^{3} \left(2240 a + 3360\right)\\
&\qquad + y \left(- 6560 a - 172640\right)\Bigr]\\
&{}+x^{16}(xz-y^2+2)\Bigl[140 a y^{5} + y^{3} \left(21840 a - 16240\right) + y \left(4480 a - 17920\right)\Bigr]\\
&{}+x^{16}(xz-y^2+2)^{2}\Bigl[y^{3} \left(4060 a - 560\right) + y \left(52080 - 112560 a\right)\Bigr]\\
&{}+x^{16}(xz-y^2+2)^{3}\Bigl[- 3080 y^{3} + y \left(- 7280 a - 3360\right)\Bigr]\\
&{}+x^{16}(xz-y^2+2)^{4}\Bigl[y^{3} \left(140 - 140 a\right) + y \left(16800 - 3360 a\right)\Bigr]\\
&{}+x^{16}(xz-y^2+2)^{5}\Bigl[y \left(1680 - 1680 a\right)\Bigr]\\
&{}+x^{16}(xz-y^2+2)^{6}\Bigl[560 a y\Bigr]\\
&{}+x^{17}\Bigl[y^{4} \left(140 - 840 a\right) + y^{2} \left(40320 a - 15680\right) + 107600 a - 103040\Bigr]\\
&{}+x^{17}(xz-y^2+2)\Bigl[1050 y^{4} + y^{2} \left(23520 a + 8400\right) - 67200 a - 15680\Bigr]\\
&{}+x^{17}(xz-y^2+2)^{2}\Bigl[y^{4} \left(70 a - 70\right) + y^{2} \left(3360 a - 12600\right) - 13440 a + 32480\Bigr]\\
&{}+x^{17}(xz-y^2+2)^{3}\Bigl[y^{2} \left(5040 a - 5040\right) - 8960 a - 21840\Bigr]\\
&{}+x^{17}(xz-y^2+2)^{4}\Bigl[- 840 a y^{2} - 8400 a - 5040\Bigr]\\
&{}+x^{17}(xz-y^2+2)^{5}\Bigl[1680 a\Bigr]\\
&{}+x^{17}(xz-y^2+2)^{6}\Bigl[560\Bigr]\\
&{}+x^{18}\Bigl[y^{5} \left(14 - 14 a\right) + y^{3} \left(2800 - 1120 a\right)\\
&\qquad + y \left(29120 a + 43680\right)\Bigr]\\
&{}+x^{18}(xz-y^2+2)\Bigl[y^{3} \left(2800 - 2800 a\right) + 2240 y\Bigr]\\
&{}+x^{18}(xz-y^2+2)^{2}\Bigl[560 a y^{3} + y \left(25200 a - 18480\right)\Bigr]\\
&{}+x^{18}(xz-y^2+2)^{3}\Bigl[1120 a y\Bigr]\\
&{}+x^{18}(xz-y^2+2)^{4}\Bigl[- 1680 y\Bigr]\\
&{}+x^{19}\Bigl[- 140 a y^{4} + y^{2} \left(7840 - 11200 a\right) - 24640 a - 38080\Bigr]\\
&{}+x^{19}(xz-y^2+2)\Bigl[- 2800 a y^{2} + 14560 a - 8960\Bigr]\\
&{}+x^{19}(xz-y^2+2)^{2}\Bigl[1680 y^{2} + 7840 a + 4480\Bigr]\\
&{}+x^{19}(xz-y^2+2)^{3}\Bigl[- 2240\Bigr]\\
&{}+x^{19}(xz-y^2+2)^{4}\Bigl[1120 a - 1120\Bigr]\\
&{}+x^{20}\Bigl[- 560 y^{3} + y \left(- 10080 a - 4480\right)\Bigr]\\
&{}+x^{20}(xz-y^2+2)^{2}\Bigl[y \left(2240 - 2240 a\right)\Bigr]\\
&{}+x^{21}\Bigl[y^{2} \left(1120 a - 1120\right) + 2240 a - 6720\Bigr]\\
&{}+x^{21}(xz-y^2+2)\Bigl[- 1120 a - 1120\Bigr]\\
&{}+x^{21}(xz-y^2+2)^{2}\Bigl[- 1120 a\Bigr]\\
&{}+x^{22}\Bigl[1120 a y\Bigr]\\
&{}+x^{23}\Bigl[448\Bigr]\\
&{}+\Bigl[x y^{2} \left(126 a - 126\right) + x \left(49 - 49 a\right) + y^{5} \left(252 a - 252\right)\\
&\qquad + y^{3} \left(700 - 700 a\right) + y \left(294 a - 322\right)\Bigr]\\
&\qquad {}\times\Bigl[(1-a)\bigl(ax(xz-y^2+2)+y\bigr)^2\\
&\qquad +\bigl((a-1)x^2-a\bigr)(xz-y^2+2)\\
&\qquad +ax(1-x^2)y-2x^4+(6-2a)x^2-2(1-a)\Bigr]\\
&{}+\Bigl[x^{2} y \left(126 - 126 a\right) + x y^{4} \left(630 - 630 a\right)\\
&\qquad + x y^{2} \left(1050 a - 1050\right) + x \left(161 - 147 a\right)\Bigr]\\
&\qquad {}\times\Bigl[(1-a)\bigl(ax(xz-y^2+2)+y\bigr)^2\\
&\qquad +\bigl((a-1)x^2-a\bigr)(xz-y^2+2)\\
&\qquad +ax(1-x^2)y-2x^4+(6-2a)x^2-2(1-a)\Bigr]^{2}\\
&{}+\Bigl[x^{3} \left(42 a - 42\right) + x^{2} y^{3} \left(840 a - 840\right)\\
&\qquad + x^{2} y \left(700 - 700 a\right)\Bigr]\\
&\qquad {}\times\Bigl[(1-a)\bigl(ax(xz-y^2+2)+y\bigr)^2\\
&\qquad +\bigl((a-1)x^2-a\bigr)(xz-y^2+2)\\
&\qquad +ax(1-x^2)y-2x^4+(6-2a)x^2-2(1-a)\Bigr]^{3}\\
&{}+\Bigl[x^{3} y^{2} \left(630 - 630 a\right) + x^{3} \left(175 a - 175\right)\Bigr]\\
&\qquad {}\times\Bigl[(1-a)\bigl(ax(xz-y^2+2)+y\bigr)^2\\
&\qquad +\bigl((a-1)x^2-a\bigr)(xz-y^2+2)\\
&\qquad +ax(1-x^2)y-2x^4+(6-2a)x^2-2(1-a)\Bigr]^{4}\\
&{}+\Bigl[x^{4} y \left(252 a - 252\right)\Bigr]\\
&\qquad {}\times\Bigl[(1-a)\bigl(ax(xz-y^2+2)+y\bigr)^2\\
&\qquad +\bigl((a-1)x^2-a\bigr)(xz-y^2+2)\\
&\qquad +ax(1-x^2)y-2x^4+(6-2a)x^2-2(1-a)\Bigr]^{5}\\
&{}+\Bigl[x^{5} \left(42 - 42 a\right)\Bigr]\\
&\qquad {}\times\Bigl[(1-a)\bigl(ax(xz-y^2+2)+y\bigr)^2\\
&\qquad +\bigl((a-1)x^2-a\bigr)(xz-y^2+2)\\
&\qquad +ax(1-x^2)y-2x^4+(6-2a)x^2-2(1-a)\Bigr]^{6}
\end{align*}

\endgroup

\clearpage

\subsubsection*{Inverse map}

The variables $(X,Y,Z)$ are independent input coordinates of the inverse.

\paragraph{Component 1.}

\begingroup\small\allowdisplaybreaks[4]

\begin{gather*}(\Phi_H^{-1})_1(X,Y,Z)=X^3-3X+Y.\end{gather*}

\endgroup

\paragraph{Component 2.}

\begingroup\footnotesize\allowdisplaybreaks[4]
\setlength{\jot}{1pt}

\begin{align*}
(\Phi_H^{-1})_2(X,Y,Z)={}&- X\\
&{}+(X^3-3X+Y)\Bigl[X^{2} - 2\Bigr]\\
&{}+(X^3-3X+Y)^{2}\Bigl[X^{7} \left(132 a - 132\right) + X^{5} \left(630 - 630 a\right) + X^{3} \left(742 a - 788\right)\\
&\qquad + X \left(181 - 120 a\right)\Bigr]\\
&{}+(X^3-3X+Y)^{3}\Bigl[- 165 X^{8} a + 1050 X^{6} a + X^{4} \left(68 - 2058 a\right) + X^{2} \left(1171 a - 196\right)\\
&\qquad - 40 a - 44\Bigr]\\
&{}+(X^3-3X+Y)^{4}\Bigl[308 X^{9} - 2040 X^{7} + X^{5} \left(162 a + 4078\right) + X^{3} \left(- 522 a - 2190\right)\\
&\qquad + X \left(196 a - 20\right)\Bigr]\\
&{}+(X^3-3X+Y)^{5}\Bigl[X^{10} \left(616 a - 616\right) + X^{8} \left(4410 - 4410 a\right)\\
&\qquad + X^{6} \left(9892 a - 10340\right) + X^{4} \left(8510 - 6810 a\right)\\
&\qquad + X^{2} \left(495 a - 1644\right) + 32 a\Bigr]\\
&{}+(X^3-3X+Y)^{6}\Bigl[- 1176 X^{11} a + 9100 X^{9} a + X^{7} \left(812 - 24444 a\right) + X^{5} \left(26474 a - 3374\right)\\
&\qquad + X^{3} \left(3082 - 9998 a\right) + X \left(384 a - 464\right)\Bigr]\\
&{}+(X^3-3X+Y)^{7}\Bigl[1764 X^{12} - 14700 X^{10} + X^{8} \left(1176 a + 42973\right)\\
&\qquad + X^{6} \left(- 5572 a - 51450\right) + X^{4} \left(7112 a + 21413\right)\\
&\qquad + X^{2} \left(224 - 2576 a\right) + 64 a - 64\Bigr]\\
&{}+\bigl(Z+(1-a)(42X^3-49X)\bigr)\Biggl\{\\
&{}(X^3-3X+Y)^{2}\Bigl[a\Bigr]\\
&{}+(X^3-3X+Y)^{3}\Bigl[2 X\Bigr]\\
&{}+(X^3-3X+Y)^{4}\Bigl[X^{2} \left(2 a - 2\right) - 5 a - 5\Bigr]\\
&{}+(X^3-3X+Y)^{5}\Bigl[- 4 X^{3} a + 10 X a\Bigr]\\
&{}+(X^3-3X+Y)^{6}\Bigl[10 X^{4} - 30 X^{2} + 4 a\Bigr]\\
&{}+(X^3-3X+Y)^{7}\Bigl[X^{5} \left(28 a - 28\right) + X^{3} \left(100 - 100 a\right) + X \left(48 a - 58\right)\Bigr]\\
&{}+(X^3-3X+Y)^{8}\Bigl[- 84 X^{6} a + 350 X^{4} a + X^{2} \left(28 - 322 a\right) + 16 a - 16\Bigr]\Biggr\}\\
&{}+\bigl(Z+(1-a)(42X^3-49X)\bigr)^{2}\Biggl\{\\
&{}(X^3-3X+Y)^{9}\Bigl[a - 1\Bigr]\Biggr\}
\end{align*}

\endgroup

\clearpage

\paragraph{Component 3.}

\begingroup\footnotesize\allowdisplaybreaks[4]

\begin{align*}
(\Phi_H^{-1})_3(X,Y,Z)={}&- 132 X^{7} a + 630 X^{5} a + X^{3} \left(46 - 788 a\right) + X \left(181 a - 62\right)\\
&{}+(X^3-3X+Y)\Bigl[X^{8} \left(429 - 264 a\right) + X^{6} \left(1260 a - 2310\right)\\
&\qquad + X^{4} \left(3562 - 1411 a\right) + X^{2} \left(29 a - 1326\right) + 44 a - 2\Bigr]\\
&{}+(X^3-3X+Y)^{2}\Bigl[X^{9} \left(902 a - 572\right) + X^{7} \left(3828 - 5928 a\right)\\
&\qquad + X^{5} \left(12184 a - 8472\right) + X^{3} \left(6618 - 7690 a\right) + X \left(511 a - 983\right)\Bigr]\\
&{}+(X^3-3X+Y)^{3}\Bigl[- 17424 X^{14} a + 166320 X^{12} a + X^{10} \left(11528 - 605878 a\right)\\
&\qquad + X^{8} \left(1047834 a - 69984\right) + X^{6} \left(139710 - 865868 a\right)\\
&\qquad + X^{4} \left(299772 a - 105335\right) + X^{2} \left(20569 - 35826 a\right) + 152 a - 208\Bigr]\\
&{}+(X^3-3X+Y)^{4}\Bigl[43560 X^{15} - 485100 X^{13} + X^{11} \left(31900 a + 2096244\right)\\
&\qquad + X^{9} \left(- 245334 a - 4385832\right) + X^{7} \left(658276 a + 4486976\right)\\
&\qquad + X^{5} \left(- 712878 a - 1953944\right) + X^{3} \left(263450 a + 219028\right)\\
&\qquad + X \left(6488 - 16752 a\right)\Bigr]\\
&{}+(X^3-3X+Y)^{5}\Bigl[X^{16} \left(108537 a - 108537\right) + X^{14} \left(1273140 - 1273140 a\right)\\
&\qquad + X^{12} \left(5868612 a - 5959804\right) + X^{10} \left(14119118 - 13362654 a\right)\\
&\qquad + X^{8} \left(15411509 a - 17656533\right) + X^{6} \left(10962294 - 8138722 a\right)\\
&\qquad + X^{4} \left(1374148 a - 2740902\right) + X^{2} \left(51300 a + 126088\right) - 2112 a - 336\Bigr]\\
&{}+(X^3-3X+Y)^{6}\Bigl[- 264264 X^{17} a + 3260400 X^{15} a + X^{13} \left(266764 - 16210456 a\right)\\
&\qquad + X^{11} \left(41388032 a - 2375388\right) + X^{9} \left(7724210 - 57080708 a\right)\\
&\qquad + X^{7} \left(40668068 a - 11053496\right) + X^{5} \left(6576582 - 12668156 a\right)\\
&\qquad + X^{3} \left(1004302 a - 1223192\right) + X \left(5408 a + 23584\right)\Bigr]\\
&{}+(X^3-3X+Y)^{7}\Bigl[608608 X^{18} - 7889700 X^{16} + X^{14} \left(653968 a + 41104248\right)\\
&\qquad + X^{12} \left(- 6220072 a - 109458596\right) + X^{10} \left(22150686 a + 156150628\right)\\
&\qquad + X^{8} \left(- 36238916 a - 113101396\right) + X^{6} \left(26789748 a + 33798166\right)\\
&\qquad + X^{4} \left(- 7506802 a - 1094458\right) + X^{2} \left(397096 a - 294864\right) + 2816\Bigr]\\
&{}+(X^3-3X+Y)^{8}\Bigl[X^{19} \left(1233232 a - 1233232\right) + X^{17} \left(16805880 - 16805880 a\right)\\
&\qquad + X^{15} \left(93184056 a - 94560256\right) + X^{13} \left(282933272 - 268948700 a\right)\\
&\qquad + X^{11} \left(428059452 a - 482380668\right) + X^{9} \left(464533414 - 364272396 a\right)\\
&\qquad + X^{7} \left(143908164 a - 232745724\right) + X^{5} \left(48510018 - 14794632 a\right)\\
&\qquad + X^{3} \left(- 1777104 a - 1872200\right) + X \left(74624 a - 30464\right)\Bigr]\\
&{}+(X^3-3X+Y)^{9}\Bigl[- 1685992 X^{20} a + 24392200 X^{18} a + X^{16} \left(2061136 - 147044758 a\right)\\
&\qquad + X^{14} \left(477435896 a - 22612296\right) + X^{12} \left(96292392 - 900050894 a\right)\\
&\qquad + X^{10} \left(987073522 a - 199641004\right) + X^{8} \left(207476496 - 592378812 a\right)\\
&\qquad + X^{6} \left(163578942 a - 100691100\right) + X^{4} \left(17775911 - 11329319 a\right)\\
&\qquad + X^{2} \left(- 183872 a - 474816\right) + 6656 a - 1536\Bigr]\\
&{}+(X^3-3X+Y)^{10}\Bigl[2535456 X^{21} - 37835840 X^{19} + X^{17} \left(3350032 a + 233477160\right)\\
&\qquad + X^{15} \left(- 38305456 a - 767587520\right) + X^{13} \left(172213356 a + 1441705300\right)\\
&\qquad + X^{11} \left(- 384911884 a - 1534382556\right) + X^{9} \left(446541820 a + 846197492\right)\\
&\qquad + X^{7} \left(- 256277268 a - 174468888\right) + X^{5} \left(60287936 a - 12454740\right)\\
&\qquad + X^{3} \left(3133312 - 2883680 a\right) + X \left(- 7680 a - 47104\right)\Bigr]\\
&{}+(X^3-3X+Y)^{11}\Bigl[X^{22} \left(3556224 a - 3556224\right) + X^{20} \left(55072080 - 55072080 a\right)\\
&\qquad + X^{18} \left(355888176 a - 360827376\right) + X^{16} \left(1300812212 - 1241692340 a\right)\\
&\qquad + X^{14} \left(2525752824 a - 2808210384\right) + X^{12} \left(3696110096 - 3009217008 a\right)\\
&\qquad + X^{10} \left(1978554914 a - 2878553048\right) + X^{8} \left(1208850804 - 589652784 a\right)\\
&\qquad + X^{6} \left(19578546 a - 215967752\right) + X^{4} \left(13022944 a + 5959328\right)\\
&\qquad + X^{2} \left(379776 - 569856 a\right) - 4096\Bigr]\\
&{}+(X^3-3X+Y)^{12}\Bigl[- 4148928 X^{23} a + 66679200 X^{21} a + X^{19} \left(5630688 - 457616880 a\right)\\
&\qquad + X^{17} \left(1749681416 a - 70284816\right) + X^{15} \left(355487160 - 4074946232 a\right)\\
&\qquad + X^{13} \left(5928291068 a - 935952332\right) + X^{11} \left(1373753220 - 5294231852 a\right)\\
&\qquad + X^{9} \left(2711832284 a - 1117135236\right) + X^{7} \left(464133572 - 672216524 a\right)\\
&\qquad + X^{5} \left(46027520 a - 78419040\right) + X^{3} \left(978688 a + 2655744\right)\\
&\qquad + X \left(10240 - 59392 a\right)\Bigr]\\
&{}+(X^3-3X+Y)^{13}\Bigl[3111696 X^{24} - 51861600 X^{22} + X^{20} \left(4148928 a + 367698744\right)\\
&\qquad + X^{18} \left(- 54232416 a - 1444921800\right) + X^{16} \left(291363408 a + 3433470817\right)\\
&\qquad + X^{14} \left(- 831187784 a - 5037568284\right) + X^{12} \left(1358704928 a + 4432878198\right)\\
&\qquad + X^{10} \left(- 1278517800 a - 2096949316\right) + X^{8} \left(651940688 a + 350528409\right)\\
&\qquad + X^{6} \left(53532864 - 150360224 a\right) + X^{4} \left(8222592 a - 10236800\right)\\
&\qquad + X^{2} \left(28672 a + 301056\right) - 4096 a\Bigr]\\
&{}+\bigl(Z+(1-a)(42X^3-49X)\bigr)\Biggl\{\\
&{}- 1\\
&{}+(X^3-3X+Y)\Bigl[- 2 X\Bigr]\\
&{}+(X^3-3X+Y)^{2}\Bigl[- 4 X^{2} + a\Bigr]\\
&{}+(X^3-3X+Y)^{3}\Bigl[- 264 X^{7} + 1260 X^{5} + X^{3} \left(- 96 a - 1472\right) + X \left(132 a + 212\right)\Bigr]\\
&{}+(X^3-3X+Y)^{4}\Bigl[X^{8} \left(198 a - 198\right) + X^{6} \left(420 - 420 a\right) + X^{4} \left(950 - 990 a\right)\\
&\qquad + X^{2} \left(1402 a - 1570\right) + 27 a + 56\Bigr]\\
&{}+(X^3-3X+Y)^{5}\Bigl[- 572 X^{9} a + 3960 X^{7} a + X^{5} \left(112 - 9240 a\right) + X^{3} \left(8000 a - 384\right)\\
&\qquad + X \left(404 - 1724 a\right)\Bigr]\\
&{}+(X^3-3X+Y)^{6}\Bigl[1716 X^{10} - 12870 X^{8} + X^{6} \left(336 a + 33684\right) + X^{4} \left(- 1360 a - 35390\right)\\
&\qquad + X^{2} \left(1782 a + 11594\right) - 408 a + 24\Bigr]\\
&{}+(X^3-3X+Y)^{7}\Bigl[X^{11} \left(5304 a - 5304\right) + X^{9} \left(42900 - 42900 a\right)\\
&\qquad + X^{7} \left(118800 a - 121968\right) + X^{5} \left(141120 - 126112 a\right)\\
&\qquad + X^{3} \left(37408 a - 55512\right) + X \left(512 a + 3016\right)\Bigr]\\
&{}+(X^3-3X+Y)^{8}\Bigl[- 16796 X^{12} a + 145860 X^{10} a + X^{8} \left(11880 - 460746 a\right)\\
&\qquad + X^{6} \left(632940 a - 62664\right) + X^{4} \left(93244 - 345978 a\right)\\
&\qquad + X^{2} \left(48088 a - 33870\right) - 32 a - 576\Bigr]\\
&{}+(X^3-3X+Y)^{9}\Bigl[54264 X^{13} - 503880 X^{11} + X^{9} \left(39468 a + 1709268\right)\\
&\qquad + X^{7} \left(- 230472 a - 2538720\right) + X^{5} \left(410256 a + 1516356\right)\\
&\qquad + X^{3} \left(- 221928 a - 226764\right) + X \left(17760 a - 6480\right)\Bigr]\\
&{}+(X^3-3X+Y)^{10}\Bigl[X^{14} \left(73752 a - 73752\right) + X^{12} \left(765660 - 765660 a\right)\\
&\qquad + X^{10} \left(2984664 a - 3045648\right) + X^{8} \left(5751450 - 5330970 a\right)\\
&\qquad + X^{6} \left(4205664 a - 5138820\right) + X^{4} \left(1824990 - 1119522 a\right)\\
&\qquad + X^{2} \left(- 32136 a - 106152\right) + 2304 a - 384\Bigr]\\
&{}+(X^3-3X+Y)^{11}\Bigl[- 123872 X^{15} a + 1333920 X^{13} a + X^{11} \left(116928 - 5562648 a\right)\\
&\qquad + X^{9} \left(11211428 a - 855288\right) + X^{7} \left(2078424 - 11045376 a\right)\\
&\qquad + X^{5} \left(4634628 a - 1865520\right) + X^{3} \left(521328 - 465580 a\right)\\
&\qquad + X \left(- 1920 a - 13056\right)\Bigr]\\
&{}+(X^3-3X+Y)^{12}\Bigl[204624 X^{16} - 2316720 X^{14} + X^{12} \left(216384 a + 10108812\right)\\
&\qquad + X^{10} \left(- 1698816 a - 21162540\right) + X^{8} \left(4610592 a + 21407148\right)\\
&\qquad + X^{6} \left(- 5046900 a - 9000432\right) + X^{4} \left(2035068 a + 664596\right)\\
&\qquad + X^{2} \left(98784 - 146304 a\right) - 1536\Bigr]\\
&{}+(X^3-3X+Y)^{13}\Bigl[X^{17} \left(296352 a - 296352\right) + X^{15} \left(3528000 - 3528000 a\right)\\
&\qquad + X^{13} \left(16547496 a - 16850904\right) + X^{11} \left(41184696 - 38661000 a\right)\\
&\qquad + X^{9} \left(46773048 a - 54298692\right) + X^{7} \left(37019472 - 27293448 a\right)\\
&\qquad + X^{5} \left(5595072 a - 10911180\right) + X^{3} \left(257472 a + 663936\right)\\
&\qquad + X \left(3840 - 22272 a\right)\Bigr]\\
&{}+(X^3-3X+Y)^{14}\Bigl[- 296352 X^{18} a + 3704400 X^{16} a + X^{14} \left(296352 - 18843048 a\right)\\
&\qquad + X^{12} \left(50007244 a - 2638944\right) + X^{10} \left(8729448 - 73544100 a\right)\\
&\qquad + X^{8} \left(58147908 a - 13293504\right) + X^{6} \left(9417912 - 21264628 a\right)\\
&\qquad + X^{4} \left(2055648 a - 2604000\right) + X^{2} \left(10752 a + 112896\right) - 2048 a\Bigr]\Biggr\}\\
&{}+\bigl(Z+(1-a)(42X^3-49X)\bigr)^{2}\Biggl\{\\
&{}(X^3-3X+Y)^{3}\Bigl[a - 1\Bigr]\\
&{}+(X^3-3X+Y)^{4}\Bigl[4 X a\Bigr]\\
&{}+(X^3-3X+Y)^{5}\Bigl[10\Bigr]\\
&{}+(X^3-3X+Y)^{7}\Bigl[- 18 a + 8\Bigr]\\
&{}+(X^3-3X+Y)^{8}\Bigl[X \left(6 - 6 a\right)\Bigr]\\
&{}+(X^3-3X+Y)^{9}\Bigl[X^{2} \left(6 a - 6\right) - 4 a - 12\Bigr]\\
&{}+(X^3-3X+Y)^{10}\Bigl[- 792 X^{7} a + 3780 X^{5} a + X^{3} \left(276 - 4728 a\right) + X \left(1086 a - 366\right)\Bigr]\\
&{}+(X^3-3X+Y)^{11}\Bigl[990 X^{8} - 6300 X^{6} + X^{4} \left(408 a + 11940\right) + X^{2} \left(- 1176 a - 5850\right)\\
&\qquad + 264 a - 24\Bigr]\\
&{}+(X^3-3X+Y)^{12}\Bigl[X^{9} \left(1848 a - 1848\right) + X^{7} \left(12240 - 12240 a\right)\\
&\qquad + X^{5} \left(24468 a - 25440\right) + X^{3} \left(16272 - 13140 a\right)\\
&\qquad + X \left(- 120 a - 1056\right)\Bigr]\\
&{}+(X^3-3X+Y)^{13}\Bigl[- 3696 X^{10} a + 26460 X^{8} a + X^{6} \left(2688 - 62040 a\right)\\
&\qquad + X^{4} \left(51060 a - 10200\right) + X^{2} \left(6894 - 9864 a\right) - 192\Bigr]\\
&{}+(X^3-3X+Y)^{14}\Bigl[7056 X^{11} - 54600 X^{9} + X^{7} \left(4872 a + 141792\right)\\
&\qquad + X^{5} \left(- 20244 a - 138600\right) + X^{3} \left(18492 a + 41496\right)\\
&\qquad + X \left(480 - 2784 a\right)\Bigr]\\
&{}+(X^3-3X+Y)^{15}\Bigl[X^{12} \left(10584 a - 10584\right) + X^{10} \left(88200 - 88200 a\right)\\
&\qquad + X^{8} \left(257838 a - 264894\right) + X^{6} \left(342132 - 308700 a\right)\\
&\qquad + X^{4} \left(128478 a - 171150\right) + X^{2} \left(1344 a + 14112\right) - 384 a\Bigr]\Biggr\}\\
&{}+\bigl(Z+(1-a)(42X^3-49X)\bigr)^{3}\Biggl\{\\
&{}(X^3-3X+Y)^{10}\Bigl[- 2\Bigr]\\
&{}+(X^3-3X+Y)^{11}\Bigl[X \left(4 a - 4\right)\Bigr]\\
&{}+(X^3-3X+Y)^{12}\Bigl[- 4 X^{2} a + 10 a\Bigr]\\
&{}+(X^3-3X+Y)^{13}\Bigl[8 X^{3} - 20 X\Bigr]\\
&{}+(X^3-3X+Y)^{14}\Bigl[X^{4} \left(20 a - 20\right) + X^{2} \left(60 - 60 a\right) - 8\Bigr]\\
&{}+(X^3-3X+Y)^{15}\Bigl[- 56 X^{5} a + 200 X^{3} a + X \left(20 - 116 a\right)\Bigr]\\
&{}+(X^3-3X+Y)^{16}\Bigl[168 X^{6} - 700 X^{4} + X^{2} \left(56 a + 588\right) - 32 a\Bigr]\Biggr\}\\
&{}+\bigl(Z+(1-a)(42X^3-49X)\bigr)^{4}\Biggl\{\\
&{}(X^3-3X+Y)^{17}\Bigl[- a\Bigr]\Biggr\}
\end{align*}

\endgroup


\clearpage
\subsection{Further transfers from the same base coordinate}

For every $N\ge1$ and $p\in K[x,y]$, the coordinate
\[
H(x,p(x,y)+x^N z,y)
\]
satisfies Theorem~\ref{thm:transfer}: modulo the nilradical of
$K[x]/(x^N)$ its centre is $-y$. The finite algorithm constructs a
companion and a polynomial inverse, retaining all powers of $x$ up to
$x^N$. This is a sufficient construction for the specified inputs; it
does not assert that every embedding of $\A^1$ into $\A^3$ admits such
a presentation.

\subsection{Computational certificate}

The file \texttt{verify.py} performs exact arithmetic in
$\Q[a,x,y,z,t]/(a^2-a+1)$ using SymPy. It checks:
\begin{enumerate}
\item the exact divisions in \eqref{eq:base-divisions} and both harmonic
      directions;
\item the five identities \eqref{eq:cert1}--\eqref{eq:cert5}, the outer
      substitutions, and the transfer identities;
\item all components of $\Phi_T(\gamma_T(t))=(t,0,0)$ and
      $\Phi_H(\gamma_H(t))=(t,0,0)$, and both inverse substitutions on
      the line $(t,0,0)$;
\item the component degrees and the relative Jacobian determinants.
\end{enumerate}
The polynomial certificates, cancellation argument, and factorization
\eqref{eq:harmonic-factorization} establish the full inverse identities.
The successful Python run is recorded in \texttt{verification.log}.
The additional script \texttt{verify\_printed\_formulas.py} reads the
actual LaTeX source and checks all twelve displayed components against
the compact construction. This separate check is essential: the original
Python certificate did not detect transcription errors in the expanded
appendices. The corrected signs are documented in the accompanying errata.

Independent exact implementations are provided in
\texttt{gap/verify\_rectification.g} and
\texttt{macaulay2/verify\_rectification.m2}. They reconstruct the
compact polynomial certificates in GAP and Macaulay2; they do not parse
the LaTeX appendices. Each implementation is accompanied by run logs
and negative controls. Tests of the Chebyshev normalization below are
also included. Detailed commands, versions and the scope of each check
are recorded in \texttt{VERIFICATION.md}.

To reproduce the Python checks:
\begin{quote}\ttfamily\small
python -m pip install -r requirements.txt\\
python -u verify.py\\
python -u verify\_printed\_formulas.py
\end{quote}
The displayed automorphisms are defined over $K\subset\C$; their
coefficients do not define a rectification by real polynomial automorphisms.

\clearpage
\subsection{Chebyshev parametrizations and the scope of the result}

Let $T_n$ denote the classical Chebyshev polynomial,
$T_n(\cos\theta)=\cos(n\theta)$, and use the normalization
$C_n(t)=2T_n(t/2)$. Thus
\[
\begin{aligned}
C_3&=t^3-3t, & C_4&=t^4-4t^2+2,\\
C_5&=t^5-5t^3+5t, & C_7&=t^7-7t^5+14t^3-7t.
\end{aligned}
\]
The real harmonic knots $H(3,4,5)$ and $H(3,5,7)$ are the trefoil and
the figure-eight, respectively; see \cite[Introduction and
Definition~1.2]{KP}. The change from $T_n$ to $C_n$ is a linear
reparametrization together with a common linear scaling of the ambient
coordinates.

The normalized Chebyshev trefoil is already covered by the construction:
\[
(C_3(t),C_4(t),C_5(t))=M(\gamma_T(t)),\qquad
M(x,y,z)=(x,y+2,z-5x).
\]
Consequently its rectifier and inverse are
\[
\Psi=\Phi_T\comp M^{-1},\qquad
\Psi^{-1}=M\comp\Phi_T^{-1},\qquad
M^{-1}(x,y,z)=(x,y-2,z+5x).
\]
In particular, $\Psi(C_3(t),C_4(t),C_5(t))=(t,0,0)$ over
$K=\Q(\sqrt{-3})$. This is a direct consequence of the explicit affine
normalization, rather than an additional hypothesis.

The harmonic model treated above is exactly $(C_3,C_5,C_7)$. The
polynomial identity
\[
C_3C_7-C_5^2+2=C_4
\]
explains the modification
$\beta(x,y,z)=(x,xz-y^2+2,y)$ connecting it to $(C_3,C_4,C_5)$.
The map $\beta$ becomes invertible after inverting $x$; the transfer
construction and its two exact divisions produce the global polynomial
automorphism.

These two cases do not establish rectifiability for all Chebyshev or
harmonic parametrizations. A further application must supply a closed
complex embedding, a coordinate restricting to its parameter, and the
hypotheses of the transfer theorem. Determining which additional
Chebyshev parametrizations satisfy these conditions remains a natural
question. In particular, the real knot classification in \cite{KP}
does not by itself establish polynomial equivalence of their complex
embeddings.


\clearpage
\begin{thebibliography}{10}
\small
\setlength{\itemsep}{7pt}

\bibitem{AM}
S.~S.~Abhyankar and T.~T.~Moh,
\emph{Embeddings of the line in the plane},
Journal f\"ur die reine und angewandte Mathematik \textbf{276}
(1975), 148--166.
\href{https://doi.org/10.1515/crll.1975.276.148}{DOI: 10.1515/crll.1975.276.148}.

\bibitem{Suzuki}
M.~Suzuki,
\emph{Propri\'et\'es topologiques des polyn\^omes de deux variables
complexes, et automorphismes alg\'ebriques de l'espace $\C^2$},
Journal of the Mathematical Society of Japan \textbf{26}(2)
(1974), 241--257.
\href{https://doi.org/10.2969/jmsj/02620241}{DOI: 10.2969/jmsj/02620241}.

\bibitem{Abhyankar}
S.~S.~Abhyankar,
\emph{On the semigroup of a meromorphic curve}, Part~I,
in \emph{Proceedings of the International Symposium on Algebraic
Geometry, Kyoto 1977}, Kinokuniya, Tokyo, 1978, 249--414.
Chapter~11, pp.~413--414.
The historical locations are documented in
\href{https://www.numdam.org/item/RSMUP_1985__74__115_0.pdf}{Craighero's
Introduction}, pp.~115--116; the original volume was not consulted here.

\bibitem{Craighero}
P.~C.~Craighero,
\emph{About Abhyankar's conjectures on space lines},
Rendiconti del Seminario Matematico della Universit\`a di Padova
\textbf{74} (1985), 115--122.
\href{https://www.numdam.org/item/RSMUP_1985__74__115_0/}{Journal archive};
\href{https://www.numdam.org/item/RSMUP_1985__74__115_0.pdf}{open PDF}.
See the Introduction and Proposition~3.

\bibitem{Shastri}
A.~R.~Shastri,
\emph{Polynomial representations of knots},
T\^ohoku Mathematical Journal (2) \textbf{44}(1) (1992), 11--17.
\href{https://doi.org/10.2748/tmj/1178227371}{DOI: 10.2748/tmj/1178227371};
\href{https://www.jstage.jst.go.jp/article/tmj1949/44/1/44_1_11/_pdf}{open PDF}.
See \S1 and \S3, Example~1.

\bibitem{KP}
P.-V.~Koseleff and D.~Pecker,
\emph{Chebyshev diagrams for two-bridge knots},
Geometriae Dedicata \textbf{150} (2011), 405--425.
\href{https://doi.org/10.1007/s10711-010-9514-7}{DOI: 10.1007/s10711-010-9514-7};
\href{https://arxiv.org/abs/0909.3281}{arXiv:0909.3281}.
See the Introduction, Definition~1.2 and the identification of $H(3,5,7)$.

\bibitem{KZ}
S.~Kaliman and M.~Zaidenberg,
\emph{Affine modifications and affine hypersurfaces with a very
transitive automorphism group},
Transformation Groups \textbf{4} (1999), 53--95.
\href{https://doi.org/10.1007/BF01236662}{DOI: 10.1007/BF01236662};
\href{https://arxiv.org/abs/math/9801076}{arXiv:math/9801076}.
See Proposition~2.1 and Corollary~2.2.

\bibitem{EV}
E.~Edo and S.~V\'en\'ereau,
\emph{Length 2 variables of $A[x,y]$ and transfer},
Annales Polonici Mathematici \textbf{76} (2001), 67--76.
\href{https://doi.org/10.4064/ap76-1-6}{DOI: 10.4064/ap76-1-6};
\href{https://www.impan.pl/shop/publication/transaction/download/product/85204}{open PDF}.
See Proposition~1, Theorem~4 and \S5.

\bibitem{EVM}
E.~Edo, A.~van den Essen and S.~Maubach,
\emph{A note on $k[z]$-automorphisms in two variables},
Journal of Pure and Applied Algebra \textbf{213} (2009), 1197--1200.
\href{https://doi.org/10.1016/j.jpaa.2008.11.040}{DOI: 10.1016/j.jpaa.2008.11.040};
\href{https://arxiv.org/abs/0809.0767}{arXiv:0809.0767}.
See \S1, Lemma~1.1, which recalls an earlier construction.

\bibitem{BD}
S.~M.~Bhatwadekar and A.~K.~Dutta,
\emph{On residual variables and stably polynomial algebras},
Communications in Algebra \textbf{21}(2) (1993), 635--645.
\href{https://doi.org/10.1080/00927879308824585}{DOI: 10.1080/00927879308824585}.
The residual-variable theorem gives the existence of a companion over
the Noetherian $\Q$-algebra $k[x]$.

\end{thebibliography}


\end{document}
